
\section{Reproduction guide: code, data, figures and numbers}\label{rep:v-repro}

Every number and figure of the Letter~\cite{v12PRL}, the paper~\cite{v12PRD} and this report is produced by the scripts of the code and data package of Ref.~\cite{DNECv12code}, whose versions share one DOI (this report corresponds to version~1.2.0, which contains it), or contained in its archived outputs, including those of the independent reimplementations.
The static-box and abelian calculations of Secs.~\ref{sec:v11-static-box} and~\ref{sec:v11-abelian} are those of the Supplemental Material of the Letter, which the package contains from version~1.2.0 (Sec.~\ref{sec:v11-repro-supplement}).
The package contains the code that produced the results, copied from the research tree with the paths rewired and the algorithms and numerical settings unchanged; the archived outputs of the original runs; a quick and a full run mode for every result; and an automatic comparison of fresh outputs with the archive under tolerances fixed in advance.
This section explains how the package is organized (Sec.~\ref{sec:v22-structure}), how to run it (Sec.~\ref{sec:v22-run}), how its results are verified (Sec.~\ref{sec:v22-verify}), where the numbers and figures come from (Secs.~\ref{sec:v22-numbers}--\ref{sec:v22-map}), and what it does not contain (Sec.~\ref{sec:v22-limits}).
The code is released under the MIT license and the data and text under CC~BY~4.0; SMDR~1.3 and the two files derived from it are under the GNU GPL, version~2 or later~\cite{Martin:2019lqd}.

\subsection{Structure}\label{sec:v22-structure}

The top level of the package is
{\small
\begin{verbatim}
run_all.py      one-command reproduction (--quick / --full)
dnec_common/    paths.py, io.py, runner.py (driver of every run.py), sm_inputs.py
casimir/ 01-03  SM Casimir weight, zero-mode energy, shape Hessian and gain
higgs/   01-12  Higgs zero-mode criterion (SM mapping ... two branches)
shape/   01-06  emptiness gate, real-moment theorem, fossil, nonlinear survival,
                Moyal and Wheeler-DeWitt checks, primordial spectra
cmb/     01-03  torus covariance, Fisher and exact KL, fossil ensemble,
                joint state
validation/     independent reimplementations: specifications, code, results,
                comparison with the main calculation
prl_supplement/ static-box and abelian-Higgs calculations of the Letter
                (its Supplemental Material, unchanged; from version 1.2.0)
figures/        style_v12.py, make_all_figures.py, one script per figure,
                data/ (inputs), out/ (PDF and PNG)
numbers/        numbers_v12.json (single source of the claim-layer numbers),
                make_numbers_tex.py, check_numbers.py
tools/          verify_expected.py, verify_summary.py, sha256sums.py
third_party/    SMDR 1.3 tarball with fetch, verify and build script
docs/technical_report/   this report (PDF and LaTeX source)
output/         run outputs (empty in the release)
\end{verbatim}}
together with \texttt{README.md}, \texttt{CHANGELOG.md}, \texttt{RELEASE\_NOTES.md}, the license, citation and environment files, \texttt{MANIFEST.json} and \texttt{SHA256SUMS.txt} (path, size and SHA-256 of every file).
Each of the 24 result directories \texttt{<chain>/NN\_<name>/} has a \texttt{README.md} (physics question, method, location in the papers, derivation number, run commands, run times, and the table of numbers that Sec.~\ref{sec:v22-map} collects), a \texttt{run.py}, the source modules, \texttt{checks.json} (comparisons and tolerances), \texttt{expected/} (archived outputs of the original runs with SHA-256 sidecars), \texttt{VERIFY.json} (results of the verification runs) and, where needed, \texttt{inputs/}.
Shared modules live in \texttt{higgs/zmlib/} and \texttt{shape/shapelib/}.
Every copied code file starts with a \texttt{PROVENANCE} line that names its original path in the research tree and the date of the copy, and every change to the code is listed in the chain's \texttt{PACKAGE\_PATCHES.json}.
The numbers D$nnn$ and C$nnn$ used throughout this report are the identifiers of the derivation files and of the calculations in the research tree; the README of each result directory gives the original paths.

\subsection{Installation and running}\label{sec:v22-run}

The package needs Python~3.12 with numpy~2.5.1, scipy~1.17.1, mpmath~1.4.1, sympy~1.14.0 and matplotlib~3.11.0 (PyMuPDF~1.27.2.2 for the figures), the versions of every recorded run:
\begin{verbatim}
conda env create -f environment.yml && conda activate dnec-v12
# or:  pip install -r requirements.txt
\end{verbatim}
Optional are psutil and threadpoolctl (for \texttt{cmb/02}), CAMB~2.0.4 (only for \texttt{cmb/01 --full}), a C compiler for SMDR~1.3 (only for the compiled refits of \texttt{higgs/01}) and Wolfram Language (for the \texttt{*.wl} cross-checks).
The package was tested on Windows~11 x64 with CPython~3.12.13; all paths go through \texttt{dnec\_common/paths.py}, so it should also run on Linux and macOS, where it has not been tested.

The entry points are
{\small
\begin{verbatim}
python run_all.py --quick              # all chains (quick) + all figures
python run_all.py --full               # lists the long computations
python run_all.py --full --yes         # runs them: about 7-8 h on 16 cores
python run_all.py --quick --chains higgs --no-figures
python higgs/run_higgs.py --quick      # one chain (also casimir, shape, cmb)
python higgs/04_cross_terms/run.py --quick    # one result directory
python figures/make_all_figures.py     # every figure from figures/data
python tools/verify_summary.py         # table of all VERIFY.json files
python numbers/check_numbers.py        # the paper numbers vs the archive
\end{verbatim}}
Every \texttt{run.py} accepts the options \texttt{--workers~N} (default: number of CPUs minus one, at most 16), \texttt{--work~DIR}, \texttt{--fresh} (discard the previous outputs of the step) and \texttt{--no-verify}.
Outputs go to \texttt{output/\allowbreak<chain>/\allowbreak<mode>/}, a tree that mirrors the archive, so that each fresh file and its archived copy share a relative path; archived upstream inputs are staged automatically when a single result directory is run.
The quick mode always recomputes something non-trivial with the packaged code (a two-dimensional node job, a DMC shard, the hazard integral, the claims layer, re-integrated nonlinear trajectories, two covariance builds) and compares it with \texttt{expected/}; the full mode repeats all original runs and can be resumed.
Multiprocessing is guarded by \texttt{if \_\_name\_\_ == "\_\_main\_\_":}, BLAS is forced to a single thread, and the code parallelizes over jobs.
The run times quoted in the READMEs refer to a 16-core Intel Core Ultra X9 378H under Windows~11.

Only \texttt{higgs/01\_sm\_mapping} uses the compiled SMDR~1.3 (one refit in the quick mode, all 38 matching-scale refits in the full mode).
The command \texttt{python third\_party/\allowbreak smdr/\allowbreak fetch\_and\_build.py all} verifies the shipped tarball by its SHA-256, unpacks, builds and tests it in about a minute; on Windows it uses MinGW-w64 gcc~5.3 with \texttt{-D\_\_USE\_MINGW\_ANSI\_STDIO=1}, a no-op \texttt{execinfo.h}, \texttt{-fcommon} and a 64~MB stack.
The stock SMDR applications are not used: \texttt{calc\_fit} crashed on the research machine and \texttt{calc\_all} did not finish within four hours, whereas the driver \texttt{c148\_fit.c}, a verbatim copy of \texttt{SMDR\_Fit\_Inputs} with the evaluation scales as arguments, reproduces the archived fits bit for bit.

\subsection{Verification}\label{sec:v22-verify}

For every result directory \texttt{checks.json} lists the compared quantities (file and key in the work tree against the archived copy, or a literal value), the location in the papers and the tolerance: for example critical top masses to $0.001\GeV$, trajectory records to $10^{-9}$ relative, and most analysis layers bit for bit.
\texttt{tools/verify\_expected.py} evaluates the checks and writes \texttt{VERIFY.json} with the differences, timings and machine; \texttt{tools/verify\_summary.py --update-readmes} copies the results into the READMEs.
Table~\ref{tab:v22-verify} gives the state of the release.
Before release the package without \texttt{output/} was copied into an empty directory, SMDR was built there from the shipped tarball, and the quick mode was run end to end; the first such self-test found a staging error in the path shims of three chains, which was corrected before release.

{\small
\setlength{\tabcolsep}{4pt}\begin{longtable}{@{}>{\raggedright\arraybackslash}p{0.27\textwidth}>{\raggedright\arraybackslash}p{0.15\textwidth}>{\raggedright\arraybackslash}p{0.12\textwidth}>{\raggedright\arraybackslash}p{0.1\textwidth}>{\raggedright\arraybackslash}p{0.15\textwidth}>{\raggedright\arraybackslash}p{0.09\textwidth}@{}}
\caption{Verification state of the code package: checks passed/total in the quick, full and reduced modes (``---'': mode not run), with the derivation and calculation numbers of each result directory and the report section that uses it. Skipped checks need optional inputs that are not shipped.}\label{tab:v22-verify}\\
\toprule Result directory & D/C no. & Report Sec. & quick & full & reduced \\ \midrule\endfirsthead
\toprule Result directory & D/C no. & Report Sec. & quick & full & reduced \\ \midrule\endhead
\bottomrule\endlastfoot
\texttt{casimir/\allowbreak{}01\_\allowbreak{}sm\_\allowbreak{}casimir\_\allowbreak{}weight} & D50 & \ref{rep:casimir} & 9/9 & 9/9 & --- \\
\texttt{casimir/\allowbreak{}02\_\allowbreak{}zero\_\allowbreak{}mode\_\allowbreak{}energy} & D54, D58, D59 & \ref{rep:zeromodes} & 9/9 & 9/9 & --- \\
\texttt{casimir/\allowbreak{}03\_\allowbreak{}shape\_\allowbreak{}hessian} & D14, D24, D55, D59 & \ref{rep:trace}, \ref{rep:shape} & 11/11 & 16/16 & --- \\
\texttt{higgs/\allowbreak{}01\_\allowbreak{}sm\_\allowbreak{}mapping} & D120, C148 & \ref{rep:h-smmap} & 14/15 (1 skipped) & 14/14 & --- \\
\texttt{higgs/\allowbreak{}02\_\allowbreak{}finite\_\allowbreak{}volume\_\allowbreak{}matching} & D110, C138 & \ref{rep:h-matching} & 9/9 & 10/10 & --- \\
\texttt{higgs/\allowbreak{}03\_\allowbreak{}zero\_\allowbreak{}mode\_\allowbreak{}2d} & D112, C140, D122, C150, T14b & \ref{rep:h-reduced}, \ref{rep:h-survival} & 10/10 & 10/10 & --- \\
\texttt{higgs/\allowbreak{}04\_\allowbreak{}cross\_\allowbreak{}terms} & D121, C149 & \ref{rep:h-cross} & 15/15 & 16/16 & --- \\
\texttt{higgs/\allowbreak{}05\_\allowbreak{}fulldim\_\allowbreak{}validation} & D123, C151, T14a & \ref{rep:v-fulldim} & 8/8 & 9/9 & --- \\
\texttt{higgs/\allowbreak{}06\_\allowbreak{}hazard\_\allowbreak{}criterion} &  & \ref{rep:h-survival}, \ref{rep:v-convergence} & 17/17 & 17/17 & --- \\
\texttt{higgs/\allowbreak{}07\_\allowbreak{}realtime\_\allowbreak{}xi} & D129, C157, C161 & \ref{rep:h-survival}, \ref{rep:h-universality} & 3/3 & 4/4 & --- \\
\texttt{higgs/\allowbreak{}08\_\allowbreak{}robustness\_\allowbreak{}w0} & D124, C152 & \ref{rep:h-survival}, \ref{rep:h-budget} & 6/6 & 7/7 & --- \\
\texttt{higgs/\allowbreak{}09\_\allowbreak{}claims} & D126, C161, D132 & \ref{rep:h-budget} & 16/16 & 14/14 & --- \\
\texttt{higgs/\allowbreak{}10\_\allowbreak{}universality} & D128, C156 & \ref{rep:h-universality} & 7/7 & 8/8 & --- \\
\texttt{higgs/\allowbreak{}11\_\allowbreak{}inflation\_\allowbreak{}stability} & D127, C155 & \ref{rep:h-universality} & 5/5 & 6/6 & --- \\
\texttt{higgs/\allowbreak{}12\_\allowbreak{}two\_\allowbreak{}branch} & D130, C159 & \ref{rep:h-survival} & 3/3 & 4/4 & --- \\
\texttt{shape/\allowbreak{}01\_\allowbreak{}emptiness\_\allowbreak{}gate} & D48, D125, D56, D59 & \ref{rep:gate} & 11/11 & 15/15 & --- \\
\texttt{shape/\allowbreak{}02\_\allowbreak{}realtime\_\allowbreak{}theorem} & D109, D113, D125 & \ref{rep:quantum-shape} & 19/19 & 16/16 & --- \\
\texttt{shape/\allowbreak{}03\_\allowbreak{}fossil} &  & \ref{rep:quantum-shape} & 16/16 & 12/12 & --- \\
\texttt{shape/\allowbreak{}04\_\allowbreak{}nonlinear\_\allowbreak{}survival} &  & \ref{rep:quantum-shape} & 13/13 & 9/9 & --- \\
\texttt{shape/\allowbreak{}05\_\allowbreak{}moyal\_\allowbreak{}wdw} &  & \ref{rep:quantum-shape} & 9/9 & 23/23 & --- \\
\texttt{shape/\allowbreak{}06\_\allowbreak{}primordial\_\allowbreak{}spectra} & D51, D30, D35, D36 & \ref{rep:spectra} & 5/5 & 5/5 & --- \\
\texttt{cmb/\allowbreak{}01\_\allowbreak{}fisher\_\allowbreak{}kl} & D52, D57 & \ref{rep:sky} & 5/5 & 3/5 (2 skipped) & --- \\
\texttt{cmb/\allowbreak{}02\_\allowbreak{}ensemble} & C145, D117 & \ref{rep:sky} & 12/12 & 13/13 & 3/3 \\
\texttt{cmb/\allowbreak{}03\_\allowbreak{}joint\_\allowbreak{}state} & C147, D119 & \ref{rep:sky} & 6/6 & 6/6 & --- \\
\midrule
\multicolumn{3}{@{}l}{total} & 238/239 & 255/257 & 3/3 \\
\end{longtable}}


The skipped check of the quick mode is the optional compiled-SMDR refit of \texttt{higgs/01}, which needs SMDR to be built; the two skipped checks of the full mode belong to \texttt{cmb/01}, whose full mode covers its C59 part only (\texttt{--full --c59-only}; the C75 part and the continuum test need inputs that are not shipped). %
Recomputing every covariance matrix of the \texttt{cmb/02} ensemble from the shipped transfer table for $\ell\le80$ (\texttt{--full}, about 70~min on 12 workers, with a 15.5~GB cache) reproduces the tabulated divergences to $2.3\times10^{-13}$, and the window edges, the peak of $P_{\rm joint}$ and its position, the integrated measure, the fossil table and the size-prior results to all quoted digits; the reduced mode of \texttt{cmb/02} exercises every code path (builds, cross-divergences, profiles, Brent searches, noise and analysis) on $\ell\le12$ against an archived reduced run. %

\subsection{Numbers}\label{sec:v22-numbers}

The claim-layer numbers of the three documents (marginal probabilities, exclusion thresholds, reach, conditions on $\xiH$, safety margins, model-error totals) are never typed by hand.
They are macros defined in \texttt{numbers.tex}, which \texttt{numbers/make\_numbers\_tex.py} generates from the single source \texttt{numbers/numbers\_v12.json}; the three LaTeX sources \texttt{\textbackslash input} the same file.
To change a number one edits only the JSON file (values and displayed precision), regenerates \texttt{numbers.tex} and copies it to the three sources (\texttt{--out PATH ...}), and reruns \texttt{python numbers/check\_numbers.py --write}, which compares every number with the archived outputs of the package.
All 51 macros are verified: each displayed value, sign included, equals the value at a named file and JSON pointer in the package rounded to the displayed precision, and none disagrees or lacks a source. %
The numbers of the static box and of the abelian model are macros of a second file, \texttt{numbers\_c170.tex}, which \texttt{prl\_supplement/\allowbreak numbers/\allowbreak make\_numbers\_c170.py} computes from the archived outputs of these calculations; the Letter and this report input the same file (Sec.~\ref{sec:v11-repro-supplement}).
Appendix~\ref{rep:app-numbers} lists the 51 macros of \texttt{numbers.tex} with their sources, together with the run05 number registry from which most of them were taken; the 52 macros of \texttt{numbers\_c170.tex} carry their sources and unrounded values in that file. %
Numbers that are not claim-layer macros are quoted in the text with a source comment in the LaTeX file; each is produced by the scripts of the package or contained in its archived outputs, including those of the independent reimplementations, and Table~\ref{tab:v22-mapping} below gives the package output and its archived copy for the numbers of each result directory.

\subsection{Figures}\label{sec:v22-figures}

\texttt{python figures/make\_all\_figures.py} (or \texttt{... <substring>} for a subset) rebuilds every data figure of the papers and of this report from copies of the archived results in \texttt{figures/data/} into \texttt{figures/out/<name>.pdf} and \texttt{.png}, with a common style file \texttt{figures/style\_v12.py}; \texttt{gen\_casimir\_slice\_shape.py} recomputes the data of the Casimir slice in about $20$~s.
Table~\ref{tab:v22-figures} lists the figures of this report, their files and the scripts that make them; schematic figures drawn in TikZ in the LaTeX source have no script.

{\small
\begin{longtable}{@{}>{\raggedright\arraybackslash}p{0.09\textwidth}>{\raggedright\arraybackslash}p{0.07\textwidth}>{\raggedright\arraybackslash}p{0.33\textwidth}>{\raggedright\arraybackslash}p{0.42\textwidth}@{}}
\caption{Figures of this report: section, image file(s) in \texttt{figures/} of the LaTeX source (copied from \texttt{figures/out/} of the package or from the paper sources), and the script of the package that makes them (in \texttt{figures/}). Schematics and diagrams drawn in TikZ have no script.}\label{tab:v22-figures}\\
\toprule Figure & Sec. & File & Script \\ \midrule\endfirsthead
\toprule Figure & Sec. & File & Script \\ \midrule\endhead
\bottomrule\endlastfoot
\ref{fig:rA-shapes} & \ref{rep:setup} & TikZ (in the LaTeX source) & -- \\
\ref{fig:rA-trace} & \ref{rep:trace} & \texttt{prd\_\allowbreak{}trace} & \texttt{fig\_\allowbreak{}prd\_\allowbreak{}trace.py} \\
\ref{fig:rA-casimir-shape} & \ref{rep:casimir} & \texttt{prd\_\allowbreak{}casimir\_\allowbreak{}shape} & \texttt{gen\_\allowbreak{}casimir\_\allowbreak{}slice\_\allowbreak{}shape.py}, \texttt{fig\_\allowbreak{}prd\_\allowbreak{}casimir\_\allowbreak{}shape.py} \\
\ref{fig:rA-holonomy-loop} & \ref{rep:zeromodes} & TikZ (in the LaTeX source) & -- \\
\ref{fig:rA-turning} & \ref{rep:turning} & TikZ (in the LaTeX source) & -- \\
\ref{fig:rb-diagrams} & \ref{rep:h-matching} & TikZ (in the LaTeX source) & -- \\
\ref{fig:rb-wavefunction} & \ref{rep:h-reduced} & \texttt{prd\_\allowbreak{}higgs\_\allowbreak{}wavefunction} & \texttt{fig\_\allowbreak{}prd\_\allowbreak{}higgs\_\allowbreak{}wavefunction.py} \\
\ref{fig:rb-gamma} & \ref{rep:h-reduced} & \texttt{prd\_\allowbreak{}gamma\_\allowbreak{}validation} & \texttt{fig\_\allowbreak{}prd\_\allowbreak{}gamma\_\allowbreak{}validation.py} \\
\ref{fig:rb-threshold} & \ref{rep:h-reduced} & \texttt{prd\_\allowbreak{}threshold} & \texttt{fig\_\allowbreak{}prd\_\allowbreak{}threshold.py} \\
\ref{fig:rb-cross} & \ref{rep:h-cross} & \texttt{prd\_\allowbreak{}cross\_\allowbreak{}terms} & \texttt{fig\_\allowbreak{}prd\_\allowbreak{}cross\_\allowbreak{}terms.py} \\
\ref{fig:rb-hazard} & \ref{rep:h-survival} & \texttt{prd\_\allowbreak{}hazard\_\allowbreak{}realtime} & \texttt{fig\_\allowbreak{}prd\_\allowbreak{}hazard\_\allowbreak{}realtime.py} \\
\ref{fig:rb-criterion} & \ref{rep:h-smmap} & \texttt{prd\_\allowbreak{}criterion} & \texttt{fig\_\allowbreak{}prd\_\allowbreak{}criterion.py} \\
\ref{fig:rb-universality} & \ref{rep:h-universality} & \texttt{prd\_\allowbreak{}universality} & \texttt{fig\_\allowbreak{}prd\_\allowbreak{}universality.py} \\
\ref{fig:rb-xi} & \ref{rep:h-universality} & \texttt{prl\_\allowbreak{}fig3\_\allowbreak{}xi} & \texttt{fig\_\allowbreak{}prl\_\allowbreak{}fig3\_\allowbreak{}xi.py} \\
\ref{fig:rA-gate} & \ref{rep:gate} & \texttt{prd\_\allowbreak{}emptiness\_\allowbreak{}gate} & \texttt{fig\_\allowbreak{}prd\_\allowbreak{}emptiness\_\allowbreak{}gate.py} \\
\ref{fig:rA-turning-state} & \ref{rep:quantum-shape} & \texttt{prd\_\allowbreak{}turning\_\allowbreak{}state} & \texttt{fig\_\allowbreak{}prd\_\allowbreak{}turning\_\allowbreak{}state.py} \\
\ref{fig:rA-fossil} & \ref{rep:quantum-shape} & \texttt{prd\_\allowbreak{}fossil\_\allowbreak{}shape} & \texttt{fig\_\allowbreak{}prd\_\allowbreak{}fossil\_\allowbreak{}shape.py} \\
\ref{fig:rA-survival} & \ref{rep:quantum-shape} & \texttt{prd\_\allowbreak{}nonlinear\_\allowbreak{}survival} & \texttt{fig\_\allowbreak{}prd\_\allowbreak{}nonlinear\_\allowbreak{}survival.py} \\
\ref{fig:rA-primordial} & \ref{rep:spectra} & \texttt{prd\_\allowbreak{}primordial} & \texttt{fig\_\allowbreak{}prd\_\allowbreak{}primordial.py} \\
\ref{fig:rA-sky} & \ref{rep:sky} & \texttt{prd\_\allowbreak{}fossil\_\allowbreak{}sky} & \texttt{fig\_\allowbreak{}prd\_\allowbreak{}fossil\_\allowbreak{}sky.py} \\
\ref{fig:v18-walls} & \ref{rep:v-fulldim} & TikZ (in the LaTeX source) & -- \\
\end{longtable}}


\subsection{From numbers to package outputs}\label{sec:v22-map}

Table~\ref{tab:v22-mapping} collects, for each result directory, the numbers quoted in the papers and in this report and the package output and key that contain them; it is generated from the tables at the end of the READMEs.
Numbers in the table are written as in the READMEs, which use some of the internal notation of the derivation files: $Q$ for the real part of the resonance energy ($\varepsilon$ in this report), $\gamma$ for the width ($\Gamma$), $u$ for the time $\tau=ht$, and $\xi$ for the zero-mode share $\fzm$ (Appendix~\ref{rep:app-notation}).
The numbers of the independent reimplementations are listed in \texttt{validation/\allowbreak COMPARISON.md} with the corresponding numbers of the main calculation, their source and the pass criterion, and \texttt{python validation/\allowbreak run.py --quick} repeats the comparison.

{\footnotesize
\begin{longtable}{@{}>{\raggedright\arraybackslash}p{0.42\textwidth}>{\raggedright\arraybackslash}p{0.53\textwidth}@{}}
\caption{Numbers of the papers and of this report and the package outputs that contain them, by result directory (from the tables at the end of each result-directory README). Output files are relative to the work tree \texttt{output/<chain>/<mode>/}; the archived copy has the same relative path below the step's \texttt{expected/} directory unless the README names another archived source; ``same'' refers to the file of the row above.}\label{tab:v22-mapping}\\
\toprule Number & Package output (work tree) : key \\ \midrule\endfirsthead
\toprule Number & Package output (work tree) : key \\ \midrule\endhead
\midrule\multicolumn{2}{r}{\emph{continued}}\\\endfoot\bottomrule\endlastfoot
\multicolumn{2}{@{}>{\raggedright\arraybackslash}p{0.97\textwidth}@{}}{\textbf{\texttt{casimir/\allowbreak{}01\_\allowbreak{}sm\_\allowbreak{}casimir\_\allowbreak{}weight}} --- signed Casimir weight of the Standard Model, spin structures and holonomies (Sec.~\ref{rep:casimir})}\\*
N\_\allowbreak{}C(AAA) = +25.515; N\_\allowbreak{}C(PPP) = \ensuremath{-}30 & \texttt{r5/\allowbreak{}C57\_\allowbreak{}spin\_\allowbreak{}structure/\allowbreak{}C57\_\allowbreak{}spin\_\allowbreak{}structures.json} : \texttt{SM\_\allowbreak{}spin\_\allowbreak{}structures/\allowbreak{}AAA/\allowbreak{}N\_\allowbreak{}C\_\allowbreak{}signed}, \texttt{.../\allowbreak{}PPP/\allowbreak{}N\_\allowbreak{}C\_\allowbreak{}signed} \\
eigenvalues +1902.24 /\allowbreak{} \ensuremath{-}893.20; \ensuremath{\kappa}\_\allowbreak{}E = 13.5286, \ensuremath{\kappa}\_\allowbreak{}T2 = \ensuremath{-}6.3524; trace identity 1e-8 & same : \texttt{SM\_\allowbreak{}spin\_\allowbreak{}structures/\allowbreak{}<s>/\allowbreak{}\{Hess\_\allowbreak{}eigs,kappa,trace\_\allowbreak{}over\_\allowbreak{}8\_\allowbreak{}3\_\allowbreak{}Z\}} \\
N\_\allowbreak{}C per spin structure; gains & same : \texttt{SM\_\allowbreak{}spin\_\allowbreak{}structures/\allowbreak{}<s>/\allowbreak{}\{N\_\allowbreak{}C\_\allowbreak{}signed,gains\_\allowbreak{}zero\_\allowbreak{}momentum\}} \\
PPP centre 11.19; PPP hypercharge max 17.79; AAA centre 22.74 /\allowbreak{} 9.86 & \texttt{r5/\allowbreak{}C68\_\allowbreak{}holonomy\_\allowbreak{}preparation/\allowbreak{}C68\_\allowbreak{}results.json} : \texttt{PPP\_\allowbreak{}centre\_\allowbreak{}111}, \texttt{PPP\_\allowbreak{}Y\_\allowbreak{}iso\_\allowbreak{}scan\_\allowbreak{}max/\allowbreak{}NC}, \texttt{AAA\_\allowbreak{}centre\_\allowbreak{}100}, \texttt{AAA\_\allowbreak{}centre\_\allowbreak{}111} \\
\addlinespace
\multicolumn{2}{@{}>{\raggedright\arraybackslash}p{0.97\textwidth}@{}}{\textbf{\texttt{casimir/\allowbreak{}02\_\allowbreak{}zero\_\allowbreak{}mode\_\allowbreak{}energy}} --- energy of the gauge and Higgs zero modes (\ensuremath{\xi} = 0.2999) and its next order (Sec.~\ref{rep:zeromodes})}\\*
\ensuremath{\xi} = 0.2993 /\allowbreak{} 0.2999 /\allowbreak{} 0.3043 & \texttt{fu/\allowbreak{}C72/\allowbreak{}C72\_\allowbreak{}summary.json} : \texttt{r\_\allowbreak{}band/\allowbreak{}one-loop mu=1/\allowbreak{}L\_\allowbreak{}b/\allowbreak{}exact\_\allowbreak{}local\_\allowbreak{}dmc\_\allowbreak{}r/\allowbreak{}lam\{0,0.01,0.1\}} \\
\ensuremath{\xi} \ensuremath{\ge} 0.2580, \ensuremath{\xi} \ensuremath{\le} 0.3083 & same : \texttt{.../\allowbreak{}lower\_\allowbreak{}operator\_\allowbreak{}bound\_\allowbreak{}r}, \texttt{.../\allowbreak{}upper\_\allowbreak{}gauss\_\allowbreak{}r/\allowbreak{}0.01} \\
DMC 2.1402(5), 4.1171(14), 12.5869(22) & same : \texttt{comparisons\_\allowbreak{}ours\_\allowbreak{}vs\_\allowbreak{}codex/\allowbreak{}*} \\
LE\_\allowbreak{}EW = 4.4478(13) /\allowbreak{} 4.4750(17); Higgs cost 1.78/\allowbreak{}L vs 0.46/\allowbreak{}L (\ensuremath{\times}3.9) & \texttt{fu/\allowbreak{}C72/\allowbreak{}C72\_\allowbreak{}dmc\_\allowbreak{}results.json} : \texttt{extrapolated/\allowbreak{}ew\_\allowbreak{}running\_\allowbreak{}lam*}; \texttt{C72\_\allowbreak{}summary.json} : \texttt{higgs\_\allowbreak{}coupled\_\allowbreak{}vs\_\allowbreak{}decoupled} \\
\ensuremath{\xi} \ensuremath{\to} 0.2826 (\ensuremath{\approx} 0.28); d\ensuremath{\xi}\_\allowbreak{}hol = \ensuremath{-}0.0140, d\ensuremath{\xi}\_\allowbreak{}Higgs = \ensuremath{-}0.0033 & \texttt{fu/\allowbreak{}C81/\allowbreak{}C81\_\allowbreak{}results.json} : \texttt{r\_\allowbreak{}first\_\allowbreak{}order/\allowbreak{}landau\_\allowbreak{}3pol}; \texttt{fu/\allowbreak{}C78/\allowbreak{}C78\_\allowbreak{}results.json} : \texttt{shifts/\allowbreak{}*} \\
\addlinespace
\multicolumn{2}{@{}>{\raggedright\arraybackslash}p{0.97\textwidth}@{}}{\textbf{\texttt{casimir/\allowbreak{}03\_\allowbreak{}shape\_\allowbreak{}hessian}} --- trace identity, no-go generalisations, gain and the next-order shape Hessian (Sec.~\ref{rep:trace}, \ref{rep:shape})}\\*
8/\allowbreak{}3 (eigenvalue at p = \ensuremath{-}2) & \texttt{fu/\allowbreak{}C73/\allowbreak{}C73\_\allowbreak{}symbolic.json} : \texttt{S2/\allowbreak{}values/\allowbreak{}-2} \\
\ensuremath{\Delta}U = (8/\allowbreak{}3)[U \ensuremath{-} Q \ensuremath{-} 3N/\allowbreak{}2] & same : \texttt{S5/\allowbreak{}*} ; \texttt{C73\_\allowbreak{}thermo.json} \\
guaranteed gain 4.500465 (NOGO-11); \ensuremath{\mathcal{G}}(13.53) = 216.46 & \texttt{fu/\allowbreak{}C73/\allowbreak{}C73\_\allowbreak{}gain.json} : \texttt{lower\_\allowbreak{}bounds/\allowbreak{}0/\allowbreak{}G via ODE}, \texttt{rows/\allowbreak{}0/\allowbreak{}gain\_\allowbreak{}closed} \\
\ensuremath{\kappa}\_\allowbreak{}E = 19.324, \ensuremath{\kappa}\_\allowbreak{}T2 = \ensuremath{-}9.074 (\ensuremath{\xi} = 0.2999) (NOGO-05/\allowbreak{}06) & \texttt{fu/\allowbreak{}C83/\allowbreak{}C83\_\allowbreak{}results.json} : \texttt{next\_\allowbreak{}order/\allowbreak{}gauss/\allowbreak{}cases/\allowbreak{}LO/\allowbreak{}\{kappa\_\allowbreak{}E,kappa\_\allowbreak{}T2\}} \\
\ensuremath{\kappa}\_\allowbreak{}E = 18.04, \ensuremath{\kappa}\_\allowbreak{}T2 = \ensuremath{-}8.34 at \ensuremath{\xi} = 0.2826 (NOGO-07); Higgs part 0.14 of the shift & same : \texttt{next\_\allowbreak{}order/\allowbreak{}gauss/\allowbreak{}cases/\allowbreak{}\{both,both\_\allowbreak{}without\_\allowbreak{}Higgs\}/\allowbreak{}kappa\_\allowbreak{}E} \\
s = 1 Laplacian \ensuremath{-}4/\allowbreak{}3; abelian zero-point ratio \ensuremath{-}8.88 & \texttt{fu/\allowbreak{}C73/\allowbreak{}C73\_\allowbreak{}zeta.json}, \texttt{C73\_\allowbreak{}abelian\_\allowbreak{}zeropoint.json} \\
\addlinespace
\multicolumn{2}{@{}>{\raggedright\arraybackslash}p{0.97\textwidth}@{}}{\textbf{\texttt{higgs/\allowbreak{}01\_\allowbreak{}sm\_\allowbreak{}mapping}} --- SM couplings at the bounce scale (D120 /\allowbreak{} C148, stage 3, run01 + run02) (Sec.~\ref{rep:h-smmap})}\\*
\ensuremath{\lambda}(\ensuremath{\mu}\_\allowbreak{}b) = \ensuremath{-}0.0097741 at (172.60, 0.1180, 125.13) & \texttt{r3/\allowbreak{}sm\_\allowbreak{}mapping/\allowbreak{}package/\allowbreak{}sm\_\allowbreak{}point.json} : \texttt{point\_\allowbreak{}172p60\_\allowbreak{}0p1180\_\allowbreak{}125p13/\allowbreak{}lam\_\allowbreak{}mu\_\allowbreak{}b} \\
y\_\allowbreak{}t(\ensuremath{\mu}\_\allowbreak{}b) = 0.4278267; g\ensuremath{_2}/\allowbreak{}g\_\allowbreak{}Y/\allowbreak{}g\ensuremath{_3} = 0.52776/\allowbreak{}0.44277/\allowbreak{}0.53565 & same file : \texttt{yt\_\allowbreak{}mu\_\allowbreak{}b}, \texttt{g2\_\allowbreak{}mu\_\allowbreak{}b}, \texttt{gY\_\allowbreak{}mu\_\allowbreak{}b}, \texttt{g3\_\allowbreak{}mu\_\allowbreak{}b} \\
\ensuremath{\mu}\_\allowbreak{}b = 3.3726\ensuremath{\times}10\ensuremath{^1}\ensuremath{^5} GeV & \texttt{sm\_\allowbreak{}point.json} : \texttt{mu\_\allowbreak{}b\_\allowbreak{}GeV} \\
\ensuremath{\lambda}\_\allowbreak{}c = \ensuremath{\lambda}(\ensuremath{\mu}\_\allowbreak{}b; m\_\allowbreak{}c) = \ensuremath{-}0.0091/\allowbreak{}\ensuremath{-}0.0121/\allowbreak{}\ensuremath{-}0.0152 & \texttt{sm\_\allowbreak{}point.json} : \texttt{lam\_\allowbreak{}mu\_\allowbreak{}b\_\allowbreak{}at\_\allowbreak{}run05\_\allowbreak{}critical\_\allowbreak{}mt} \\
\ensuremath{\bar{m}}\_\allowbreak{}t(\ensuremath{\bar{m}}\_\allowbreak{}t) 162.7281/\allowbreak{}163.1637/\allowbreak{}163.6012 (\ensuremath{\alpha}\_\allowbreak{}s 0.1180; also 0.1171/\allowbreak{}0.1189) & \texttt{r3/\allowbreak{}sm\_\allowbreak{}mapping/\allowbreak{}run02/\allowbreak{}R3\_\allowbreak{}mbar\_\allowbreak{}run05.json} : \texttt{new\_\allowbreak{}central/\allowbreak{}<\ensuremath{\alpha}\_\allowbreak{}s>/\allowbreak{}P1=*/\allowbreak{}mbar\_\allowbreak{}mbar\_\allowbreak{}QCD4L} \\
compiled SMDR \ensuremath{-} Python at Q = M\_\allowbreak{}t; \ensuremath{\lambda}(\ensuremath{\mu}\_\allowbreak{}b) SMDR-RG /\allowbreak{} own RG /\allowbreak{} Python & \texttt{r3/\allowbreak{}sm\_\allowbreak{}mapping/\allowbreak{}run02/\allowbreak{}C148r2\_\allowbreak{}summary.json} : \texttt{A\_\allowbreak{}match\_\allowbreak{}at\_\allowbreak{}Mt\_\allowbreak{}compiled\_\allowbreak{}minus\_\allowbreak{}python}, \texttt{B\_\allowbreak{}mu\_\allowbreak{}b} \\
compiled matching-scale variation +0.0331/\allowbreak{}\ensuremath{-}0.0186 GeV & \texttt{C148r2\_\allowbreak{}summary.json} : \texttt{E\_\allowbreak{}budget\_\allowbreak{}P1=0.5/\allowbreak{}a\_\allowbreak{}matching\_\allowbreak{}scale\_\allowbreak{}compiledSMDR/\allowbreak{}common\_\allowbreak{}xi} \\
benchmarks (BKPV critical M\_\allowbreak{}t 171.453; r2\ensuremath{\to}r3 chain +0.063 GeV) & \texttt{r3/\allowbreak{}sm\_\allowbreak{}mapping/\allowbreak{}run01/\allowbreak{}C148\_\allowbreak{}bench.json} (full mode) \\
absolute-stability bound m\_\allowbreak{}stab = 170.99 GeV (m0 \ensuremath{-} m\_\allowbreak{}stab = 0.10 GeV) & \texttt{w/\allowbreak{}abs\_\allowbreak{}stability/\allowbreak{}abs\_\allowbreak{}stability.json} : \texttt{central/\allowbreak{}m\_\allowbreak{}stab\_\allowbreak{}pole}, \texttt{central/\allowbreak{}m0\_\allowbreak{}pole} \\
\addlinespace
\multicolumn{2}{@{}>{\raggedright\arraybackslash}p{0.97\textwidth}@{}}{\textbf{\texttt{higgs/\allowbreak{}02\_\allowbreak{}finite\_\allowbreak{}volume\_\allowbreak{}matching}} --- consistent one-loop finite-volume matching (D110 /\allowbreak{} C138, stage 2) (Sec.~\ref{rep:h-matching})}\\*
T\_\allowbreak{}B/\allowbreak{}2 = \ensuremath{-}0.1128925, c\ensuremath{_2} = 0.0062594 & \texttt{r2/\allowbreak{}higgs\_\allowbreak{}matching/\allowbreak{}run01/\allowbreak{}C138\_\allowbreak{}checks.json} : \texttt{R14\_\allowbreak{}fv\_\allowbreak{}internal/\allowbreak{}TB\_\allowbreak{}half}, \texttt{c2P} \\
\ensuremath{\gamma}\_\allowbreak{}\ensuremath{\phi}\textasciicircum{}C(\ensuremath{\mu}\_\allowbreak{}b) = +2.2\ensuremath{\times}10\ensuremath{^-}\ensuremath{^4} & \texttt{C138\_\allowbreak{}checks.json} : \texttt{RG\_\allowbreak{}one\_\allowbreak{}loop\_\allowbreak{}consistency\_\allowbreak{}lam-0.010/\allowbreak{}gamma\_\allowbreak{}phi\_\allowbreak{}Coulomb} \\
\ensuremath{\gamma}(\ensuremath{\lambda} = \ensuremath{-}0.010, v11) MC+CW 9.871\ensuremath{\times}10\ensuremath{^-}\ensuremath{^4}; MU 2.051\ensuremath{\times}10\ensuremath{^-}\ensuremath{^3} & \texttt{C138\_\allowbreak{}static.json} : \texttt{table/\allowbreak{}v11/\allowbreak{}MC/\allowbreak{}-0.01/\allowbreak{}gamma}, \texttt{table/\allowbreak{}v11/\allowbreak{}MU/\allowbreak{}-0.01/\allowbreak{}gamma} \\
MC+CW (D110) 172.349/\allowbreak{}172.768 (R14 mapping) & \texttt{C138\_\allowbreak{}summary.json} : \texttt{["central\_\allowbreak{}MC+CW\_\allowbreak{}mu=1/\allowbreak{}L"]} (full) \\
\ensuremath{\mu} \ensuremath{\in} [\ensuremath{\tfrac12}, 2\ensuremath{\pi}]/\allowbreak{}L items & \texttt{C138\_\allowbreak{}summary.json} : \texttt{error\_\allowbreak{}budget\_\allowbreak{}GeV} (full) \\
matching vs zero-point confinement: local slope lowered by 10--33\% at r = 1--5; barrier above the level U\_\allowbreak{}ad 1.565 \ensuremath{\to} 0.715 (\ensuremath{\lambda} = \ensuremath{-}0.010), 1.211 \ensuremath{\to} 0.534 (\ensuremath{-}0.013) & \texttt{w/\allowbreak{}matching\_\allowbreak{}potential/\allowbreak{}matching\_\allowbreak{}potential.json} : \texttt{headline}, \texttt{lambda/\allowbreak{}<\ensuremath{\lambda}>/\allowbreak{}slope}, \texttt{lambda/\allowbreak{}<\ensuremath{\lambda}>/\allowbreak{}barrier\_\allowbreak{}above\_\allowbreak{}Q} (\texttt{matching\_\allowbreak{}potential.py}) \\
\addlinespace
\multicolumn{2}{@{}>{\raggedright\arraybackslash}p{0.97\textwidth}@{}}{\textbf{\texttt{higgs/\allowbreak{}03\_\allowbreak{}zero\_\allowbreak{}mode\_\allowbreak{}2d}} --- 2D (r, \textbar{}w\textbar{}) zero-mode model and the direct 2D hazard (D112 /\allowbreak{} C140; D122 /\allowbreak{} C150 = T14b) (Sec.~\ref{rep:h-reduced}, \ref{rep:h-survival})}\\*
\ensuremath{\lambda} = 0: 2D\_\allowbreak{}ex 4.45086, 2D\_\allowbreak{}G 4.5077 vs DMC 4.4478(13) & \texttt{r2/\allowbreak{}higgs\_\allowbreak{}beyond\_\allowbreak{}bo/\allowbreak{}run01/\allowbreak{}C140\_\allowbreak{}lam0.json} : \texttt{Q}, \texttt{test/\allowbreak{}Q\_\allowbreak{}2D\_\allowbreak{}ex}, \texttt{test/\allowbreak{}minus\_\allowbreak{}C72} \\
172.5185/\allowbreak{}173.0140/\allowbreak{}173.5155 & \texttt{r3/\allowbreak{}twod\_\allowbreak{}direct\_\allowbreak{}sm/\allowbreak{}run01/\allowbreak{}C150\_\allowbreak{}summary.json} : \texttt{main\_\allowbreak{}2D\_\allowbreak{}direct\_\allowbreak{}as0.1180/\allowbreak{}critical\_\allowbreak{}mt} \\
172.4671/\allowbreak{}172.9412 & \texttt{\ldots{}/\allowbreak{}run01/\allowbreak{}C150\_\allowbreak{}summary.json} : \texttt{twoD\_\allowbreak{}G\_\allowbreak{}direct/\allowbreak{}critical\_\allowbreak{}mt} \\
MU 172.2099/\allowbreak{}172.5837 \ldots{} & \texttt{\ldots{}/\allowbreak{}run01/\allowbreak{}C150\_\allowbreak{}summary.json} : \texttt{convergence\_\allowbreak{}table} \\
172.5826/\allowbreak{}173.0775/\allowbreak{}173.5784 & \texttt{r3/\allowbreak{}twod\_\allowbreak{}direct\_\allowbreak{}sm/\allowbreak{}run02/\allowbreak{}C150\_\allowbreak{}summary.json} : \texttt{main\_\allowbreak{}2D\_\allowbreak{}direct\_\allowbreak{}as0.1180/\allowbreak{}critical\_\allowbreak{}mt} \\
\ensuremath{\gamma}(m\_\allowbreak{}t, u) at 11 nodes & \texttt{r3/\allowbreak{}twod\_\allowbreak{}direct\_\allowbreak{}sm/\allowbreak{}run02/\allowbreak{}shards2d/\allowbreak{}2Dex\_\allowbreak{}*.json} : \texttt{cols/\allowbreak{}*/\allowbreak{}gamma2} \\
\addlinespace
\multicolumn{2}{@{}>{\raggedright\arraybackslash}p{0.97\textwidth}@{}}{\textbf{\texttt{higgs/\allowbreak{}04\_\allowbreak{}cross\_\allowbreak{}terms}} --- zero-mode /\allowbreak{} non-zero-mode cross terms and the run05 node widths (D121 /\allowbreak{} C149, run02 \ensuremath{\to} run05) (Sec.~\ref{rep:h-cross})}\\*
m\_\allowbreak{}c\textasciicircum{}pole 172.5024/\allowbreak{}172.9595/\allowbreak{}173.4184 (\ensuremath{\alpha}\_\allowbreak{}s 0.1180) & \texttt{r3/\allowbreak{}holonomy\_\allowbreak{}cross/\allowbreak{}run05/\allowbreak{}C149\_\allowbreak{}run05\_\allowbreak{}summary.json} : \texttt{crit/\allowbreak{}0.1180/\allowbreak{}P1=*} \\
\ensuremath{\alpha}\_\allowbreak{}s 0.1171: 172.156/\allowbreak{}172.611/\allowbreak{}173.069; 0.1189: 172.849/\allowbreak{}173.308/\allowbreak{}173.768 & same : \texttt{crit/\allowbreak{}0.1171}, \texttt{crit/\allowbreak{}0.1189} \\
P\ensuremath{_1}(172.60) = 0.8426; P\ensuremath{_1}(172.52) = 0.8909 & same : \texttt{P1/\allowbreak{}172.6}, \texttt{P1/\allowbreak{}172.52} \\
P\ensuremath{_1} grid 172.0\ldots{}173.8 & same : \texttt{P1\_\allowbreak{}grid} \\
\ensuremath{\lambda}\_\allowbreak{}c(\ensuremath{\mu}\_\allowbreak{}b) \ensuremath{-}0.0091339/\allowbreak{}\ensuremath{-}0.0121415/\allowbreak{}\ensuremath{-}0.0151849 & same : \texttt{lam\_\allowbreak{}c} \\
\ensuremath{-}0.0802/\allowbreak{}\ensuremath{-}0.1180/\allowbreak{}\ensuremath{-}0.1600 vs no cross term & same : \texttt{shift\_\allowbreak{}vs\_\allowbreak{}T14b/\allowbreak{}0.1180} \\
RG closure 0.994/\allowbreak{}0.973/\allowbreak{}0.993/\allowbreak{}0.990 & \texttt{r3/\allowbreak{}holonomy\_\allowbreak{}cross/\allowbreak{}run05/\allowbreak{}C149\_\allowbreak{}coulomb\_\allowbreak{}check.json} : \texttt{<r>/\allowbreak{}closure\_\allowbreak{}Z} \\
top chiral G(m\_\allowbreak{}t\ensuremath{^2}), B/\allowbreak{}(q\_\allowbreak{}\ensuremath{\Phi}\ensuremath{^2} m\ensuremath{^2}) & \texttt{r3/\allowbreak{}holonomy\_\allowbreak{}cross/\allowbreak{}run04/\allowbreak{}C149\_\allowbreak{}top\_\allowbreak{}chiral.json} \\
node widths \ensuremath{\gamma}(m\_\allowbreak{}t, u) with cross terms & \texttt{r3/\allowbreak{}holonomy\_\allowbreak{}cross/\allowbreak{}run05/\allowbreak{}shards/\allowbreak{}X\_\allowbreak{}\{main,als\}\_\allowbreak{}u*.json} : \texttt{cols/\allowbreak{}*/\allowbreak{}gamma} \\
split by zero-mode field at P\ensuremath{_1} = 0.5: SU(3)--top \ensuremath{-}0.0403, vev mixing + Coulomb +0.0275, top chiral \ensuremath{-}0.0219 (R5-049\ldots{}053) & \texttt{w/\allowbreak{}claims\_\allowbreak{}run05/\allowbreak{}run01/\allowbreak{}cross\_\allowbreak{}decomp/\allowbreak{}cross\_\allowbreak{}decomp\_\allowbreak{}summary.json} (\texttt{cross\_\allowbreak{}decomp.py post}, from the 66 archived node shards) \\
\addlinespace
\multicolumn{2}{@{}>{\raggedright\arraybackslash}p{0.97\textwidth}@{}}{\textbf{\texttt{higgs/\allowbreak{}05\_\allowbreak{}fulldim\_\allowbreak{}validation}} --- full-dimensional check of the reduced zero-mode model (D123 /\allowbreak{} C151 /\allowbreak{} T14a) (Sec.~\ref{rep:v-fulldim})}\\*
\ensuremath{\Delta}Q (same wall) \ensuremath{-}0.0032(8) (\ensuremath{\lambda} = \ensuremath{-}0.010), \ensuremath{-}0.0034(9) (\ensuremath{-}0.013) & \texttt{r3/\allowbreak{}fulldim\_\allowbreak{}dmc/\allowbreak{}run01/\allowbreak{}C151\_\allowbreak{}summary.json} : \texttt{lam\_\allowbreak{}neg/\allowbreak{}-0.01/\allowbreak{}dQ\_\allowbreak{}ex\_\allowbreak{}main}, \texttt{lam\_\allowbreak{}neg/\allowbreak{}-0.013/\allowbreak{}dQ\_\allowbreak{}ex\_\allowbreak{}main} \\
\ensuremath{\lambda} = 0: Q\_\allowbreak{}full 4.44659(94) vs C72 4.4478(13) & same : \texttt{lam0\_\allowbreak{}regression} \\
tail amplitude A = 0.92(16)/\allowbreak{}1.03(7); WKB 0.937/\allowbreak{}0.973; CAP d ln \ensuremath{\gamma}/\allowbreak{}dQ & same : \texttt{lam\_\allowbreak{}neg/\allowbreak{}*/\allowbreak{}tail\_\allowbreak{}fit}, \texttt{wkb}, \texttt{conversion} \\
3D collective coordinate \ensuremath{\gamma}\_\allowbreak{}3D/\allowbreak{}\ensuremath{\gamma}\_\allowbreak{}2D 0.9992--1.0000 & same : \texttt{model\_\allowbreak{}3D\_\allowbreak{}zeta} \\
\ensuremath{\Delta}m\_\allowbreak{}c +0.0043(10)/\allowbreak{}+0.0070(17) GeV & same : \texttt{delta\_\allowbreak{}m\_\allowbreak{}c} \\
DMC energies per wall and tau & \texttt{r3/\allowbreak{}fulldim\_\allowbreak{}dmc/\allowbreak{}run01/\allowbreak{}T14a\_\allowbreak{}dmc\_\allowbreak{}combined.json}, \texttt{shards/\allowbreak{}*.json} \\
\addlinespace
\multicolumn{2}{@{}>{\raggedright\arraybackslash}p{0.97\textwidth}@{}}{\textbf{\texttt{higgs/\allowbreak{}06\_\allowbreak{}hazard\_\allowbreak{}criterion}} --- hazard along the turn, P\ensuremath{_1}(m\_\allowbreak{}t, \ensuremath{\alpha}\_\allowbreak{}s), critical masses, \ensuremath{\lambda}\_\allowbreak{}c, convergence table (Sec.~\ref{rep:h-survival}, \ref{rep:v-convergence})}\\*
m\_\allowbreak{}c\textasciicircum{}pole 172.5024/\allowbreak{}172.9595/\allowbreak{}173.4184 & \texttt{r3/\allowbreak{}holonomy\_\allowbreak{}cross/\allowbreak{}run05/\allowbreak{}C149\_\allowbreak{}run05\_\allowbreak{}summary.json} : \texttt{crit/\allowbreak{}0.1180} \\
m\_\allowbreak{}c at \ensuremath{\alpha}\_\allowbreak{}s 0.1171/\allowbreak{}0.1189; slope dm\_\allowbreak{}c/\allowbreak{}d\ensuremath{\alpha}\_\allowbreak{}s \ensuremath{\approx} 387 GeV & same : \texttt{crit} \\
\ensuremath{\lambda}\_\allowbreak{}c \ensuremath{-}0.0091/\allowbreak{}\ensuremath{-}0.0121/\allowbreak{}\ensuremath{-}0.0152 & same : \texttt{lam\_\allowbreak{}c} \\
P\ensuremath{_1}(172.60) 0.8426; P\ensuremath{_1}(172.52) 0.8909; P\ensuremath{_1}(172.2) 0.9850; P\ensuremath{_1}(173.0) 0.4549 & same : \texttt{P1} \\
P\ensuremath{_1}(m\_\allowbreak{}t) grid and I(m\_\allowbreak{}t) & \texttt{r3/\allowbreak{}hazard\_\allowbreak{}criterion/\allowbreak{}hazard\_\allowbreak{}table.json} : \texttt{P1/\allowbreak{}0.1180}, \texttt{I/\allowbreak{}0.1180} \\
\ensuremath{\bar{m}}\_\allowbreak{}t 162.7281/\allowbreak{}163.1637/\allowbreak{}163.6012 & \texttt{r3/\allowbreak{}sm\_\allowbreak{}mapping/\allowbreak{}run02/\allowbreak{}R3\_\allowbreak{}mbar\_\allowbreak{}run05.json} : \texttt{new\_\allowbreak{}central} \\
MU \ensuremath{\to} \ldots{} \ensuremath{\to} 2D direct (R14) \ensuremath{\to} +NNLO \ensuremath{\to} +cross (run02) \ensuremath{\to} run04 \ensuremath{\to} run05 & \texttt{r3/\allowbreak{}hazard\_\allowbreak{}criterion/\allowbreak{}hazard\_\allowbreak{}table.json} : \texttt{convergence\_\allowbreak{}table\_\allowbreak{}0p1180} \\
quadrature /\allowbreak{} interpolation sensitivity & \texttt{hazard\_\allowbreak{}table.json} : \texttt{max\_\allowbreak{}abs\_\allowbreak{}diff\_\allowbreak{}*} \\
\addlinespace
\multicolumn{2}{@{}>{\raggedright\arraybackslash}p{0.97\textwidth}@{}}{\textbf{\texttt{higgs/\allowbreak{}07\_\allowbreak{}realtime\_\allowbreak{}xi}} --- real-time evolution through the turn and the condition on \ensuremath{\xi}\_\allowbreak{}H (D129 /\allowbreak{} C157; C161 run05 version) (Sec.~\ref{rep:h-survival}, \ref{rep:h-universality})}\\*
\ensuremath{\xi}\_\allowbreak{}c(172.6) at run05 = 0.0738 \ensuremath{\pm} 0.0025 (run02: 0.0741); P\ensuremath{_1} = 0.5 point 0.1205 & \texttt{w/\allowbreak{}claims\_\allowbreak{}run05/\allowbreak{}run01/\allowbreak{}xi\_\allowbreak{}run05\_\allowbreak{}summary.json} \\
\ensuremath{\xi}\_\allowbreak{}c(m\_\allowbreak{}t) table, EW image & \texttt{w/\allowbreak{}claims\_\allowbreak{}run05/\allowbreak{}run01/\allowbreak{}xi\_\allowbreak{}run05/\allowbreak{}C157\_\allowbreak{}xi\_\allowbreak{}quick.json} \\
universal \ensuremath{\xi}\_\allowbreak{}c(P\ensuremath{_1}): residuals of 16 run02 variants and run05 points & \texttt{xi\_\allowbreak{}run05\_\allowbreak{}summary.json}; \texttt{xi\_\allowbreak{}run05\_\allowbreak{}compare.csv} \\
SM band 0.0750 \ensuremath{\pm} 0.0130, \ensuremath{\alpha}\_\allowbreak{}s band 0.0795 \ensuremath{\pm} 0.0397 (m\_\allowbreak{}t = 172.6) & archive only \\
2D vs 1D: \ensuremath{\Delta}P\_\allowbreak{}surv \ensuremath{\le} 2\ensuremath{\times}10\ensuremath{^-}\ensuremath{^4} & archive only \\
\addlinespace
\multicolumn{2}{@{}>{\raggedright\arraybackslash}p{0.97\textwidth}@{}}{\textbf{\texttt{higgs/\allowbreak{}08\_\allowbreak{}robustness\_\allowbreak{}w0}} --- background sensitivity, inflation-scale item, B3-rule w\ensuremath{_0} (D124 /\allowbreak{} C152) (Sec.~\ref{rep:h-survival}, \ref{rep:h-budget})}\\*
\ensuremath{\pm}0.035/\allowbreak{}\ensuremath{\pm}0.067/\allowbreak{}\ensuremath{\pm}0.112 GeV (r \ensuremath{\in} [10\ensuremath{^-}\ensuremath{^3}, 10\ensuremath{^-}\ensuremath{^2}]); r = 10\ensuremath{^-}\ensuremath{^4} tail \ensuremath{-}0.089/\allowbreak{}\ensuremath{-}0.1725/\allowbreak{}\ensuremath{-}0.2795 & \texttt{r3/\allowbreak{}higgs\_\allowbreak{}robustness/\allowbreak{}run01/\allowbreak{}C152\_\allowbreak{}input\_\allowbreak{}dmc.json} : \texttt{vi\_\allowbreak{}inflation\_\allowbreak{}scale} \\
d ln I/\allowbreak{}d\ensuremath{\lambda}\_\allowbreak{}b \ensuremath{-}805.4/\allowbreak{}\ensuremath{-}487.8/\allowbreak{}\ensuremath{-}319.2 & same : \texttt{slope\_\allowbreak{}update/\allowbreak{}dlnI\_\allowbreak{}dlam\_\allowbreak{}b} \\
e-fold finite differences of hL\_\allowbreak{}b (\ensuremath{-}0.170/\allowbreak{}\ensuremath{-}0.270/\allowbreak{}\ensuremath{-}0.394) & \texttt{C152\_\allowbreak{}post.json} \\
w\ensuremath{_0}(B3) 0.9945/\allowbreak{}0.9907/\allowbreak{}0.9852/\allowbreak{}0.9801; critical 172.4819/\allowbreak{}172.9473/\allowbreak{}173.4076 (run02 base) & \texttt{C152\_\allowbreak{}mc\_\allowbreak{}w0.json} : \texttt{with\_\allowbreak{}B3\_\allowbreak{}w0\_\allowbreak{}of\_\allowbreak{}mt} ; \texttt{C152\_\allowbreak{}w0\_\allowbreak{}B3.json} \\
h/\allowbreak{}M\_\allowbreak{}P 6.054e-6 (Higgs calculation) vs 6.116e-6 (instant reheating): +0.0006/\allowbreak{}+0.0012/\allowbreak{}+0.0020 GeV & \texttt{C152\_\allowbreak{}accounting.json} \\
\addlinespace
\multicolumn{2}{@{}>{\raggedright\arraybackslash}p{0.97\textwidth}@{}}{\textbf{\texttt{higgs/\allowbreak{}09\_\allowbreak{}claims}} --- claims layer at the run05 centre (D126 v2 + correction addendum; C161 /\allowbreak{} D132) (Sec.~\ref{rep:h-budget})}\\*
\ensuremath{\bar{P}}\ensuremath{_1} direct 0.7302 (run02: 0.7271) & \texttt{w/\allowbreak{}claims\_\allowbreak{}run05/\allowbreak{}run01/\allowbreak{}C154\_\allowbreak{}correction\_\allowbreak{}summary.json} : \texttt{marginal/\allowbreak{}equal\_\allowbreak{}mass/\allowbreak{}D1\_\allowbreak{}direct\_\allowbreak{}pole/\allowbreak{}Pbar} \\
\ensuremath{\bar{P}}\ensuremath{_1} for \ensuremath{\sigma}-pole, MSR, \ensuremath{\bar{m}}\_\allowbreak{}t, ATLAS+CMS, CMS, BKPV/\allowbreak{}FJ variant & same : \texttt{marginal/\allowbreak{}equal\_\allowbreak{}mass/\allowbreak{}*} \\
Bayes B\_\allowbreak{}ex, B\_\allowbreak{}sel & same : \texttt{bayes} \\
HL-LHC separation 1.7--2.5\ensuremath{\sigma} & same : \texttt{HLLHC} \\
critical masses incl. P\ensuremath{_1} = 0.05; exclusion thresholds; FCC-ee 5.3\ensuremath{\sigma} & \texttt{w/\allowbreak{}claims\_\allowbreak{}run05/\allowbreak{}run01/\allowbreak{}C154\_\allowbreak{}summary.json} : \texttt{critical\_\allowbreak{}mt}, \texttt{exclusion\_\allowbreak{}mt}, \texttt{future}; \texttt{future\_\allowbreak{}reach.csv}, \texttt{exclusion\_\allowbreak{}mt.csv} \\
P\ensuremath{_1}(m\_\allowbreak{}t, \ensuremath{\alpha}\_\allowbreak{}s) grid, contours, ellipses, \ensuremath{\xi} curves & \texttt{fig2a\_\allowbreak{}*.csv}, \texttt{fig2b\_\allowbreak{}*} in \texttt{w/\allowbreak{}claims\_\allowbreak{}run05/\allowbreak{}run01/\allowbreak{}} \\
model composite +0.0721/\allowbreak{}\ensuremath{-}0.0842 at P\ensuremath{_1} = 0.5; matching row +0.012/\allowbreak{}\ensuremath{-}0.049 & \texttt{w/\allowbreak{}claims\_\allowbreak{}run05/\allowbreak{}run01/\allowbreak{}table1\_\allowbreak{}rebuild.json}, \texttt{table1\_\allowbreak{}error\_\allowbreak{}budget\_\allowbreak{}run05.csv} \\
P\ensuremath{_2} critical 172.380/\allowbreak{}172.764/\allowbreak{}173.131; B3 172.486/\allowbreak{}172.953/\allowbreak{}173.414 & \texttt{w/\allowbreak{}claims\_\allowbreak{}run05/\allowbreak{}run01/\allowbreak{}crit\_\allowbreak{}variants.json} : \texttt{run05} \\
independent MC \ensuremath{\bar{P}}\ensuremath{_1} (0.7304 direct) & \texttt{w/\allowbreak{}claims\_\allowbreak{}run05/\allowbreak{}run01/\allowbreak{}pbar\_\allowbreak{}direct\_\allowbreak{}mc.json} : \texttt{results} \\
FCC-ee exclusion strength at \ensuremath{\hat m} = 173.68 GeV: central P\ensuremath{_1} = 0.023, B\_\allowbreak{}ex = 0.033 (\texttt{\textbackslash PoneAtExclFCC}, \texttt{\textbackslash BexFCC}) & \texttt{w/\allowbreak{}claims\_\allowbreak{}run05/\allowbreak{}run01/\allowbreak{}fcc\_\allowbreak{}exclusion.json} : \texttt{headline} (\texttt{fcc\_\allowbreak{}exclusion.py}) \\
\addlinespace
\multicolumn{2}{@{}>{\raggedright\arraybackslash}p{0.97\textwidth}@{}}{\textbf{\texttt{higgs/\allowbreak{}10\_\allowbreak{}universality}} --- how universal is \ensuremath{\lambda}\_\allowbreak{}c? (D128 /\allowbreak{} C156) (Sec.~\ref{rep:h-universality})}\\*
b\_\allowbreak{}eff 1.3992/\allowbreak{}1.2991/\allowbreak{}1.1893; d\ensuremath{\lambda}\_\allowbreak{}c/\allowbreak{}d ln \ensuremath{\eta} 0.00124/\allowbreak{}0.00206/\allowbreak{}0.00315; dm\_\allowbreak{}c/\allowbreak{}d ln \ensuremath{\eta} \ensuremath{-}0.1898/\allowbreak{}\ensuremath{-}0.3113/\allowbreak{}\ensuremath{-}0.4732 & \texttt{r3/\allowbreak{}universality/\allowbreak{}run01/\allowbreak{}C156\_\allowbreak{}summary.json} : \texttt{eta} \\
\ensuremath{\partial}ln\textbar{}\ensuremath{\lambda}\_\allowbreak{}c\textbar{}/\allowbreak{}\ensuremath{\partial}ln(g\ensuremath{_2}, g\_\allowbreak{}Y, y\_\allowbreak{}t) 0.663/\allowbreak{}0.163/\allowbreak{}0.627 & same : \texttt{couplings} \\
N\_\allowbreak{}C 13.124/\allowbreak{}20.871 \ensuremath{\to} 21.068/\allowbreak{}23.669 & \texttt{C156\_\allowbreak{}centre.json} \\
flux sector 100.719/\allowbreak{}L, ratio 2.3565 & \texttt{C156\_\allowbreak{}twist.json} \\
run02 repair & \texttt{r3/\allowbreak{}universality/\allowbreak{}run02/\allowbreak{}C156\_\allowbreak{}run02.json} \\
\addlinespace
\multicolumn{2}{@{}>{\raggedright\arraybackslash}p{0.97\textwidth}@{}}{\textbf{\texttt{higgs/\allowbreak{}11\_\allowbreak{}inflation\_\allowbreak{}stability}} --- all-mode Higgs stability during inflation (D127 /\allowbreak{} C155) (Sec.~\ref{rep:h-universality})}\\*
\ensuremath{\xi}\_\allowbreak{}inf(EW) 0.04877/\allowbreak{}0.05512/\allowbreak{}0.06044 (\ensuremath{\alpha}\_\allowbreak{}s 0.118) & \texttt{r3/\allowbreak{}inflation\_\allowbreak{}stability/\allowbreak{}run01/\allowbreak{}xi\_\allowbreak{}inf\_\allowbreak{}table.csv} : column \texttt{xi\_\allowbreak{}inf\_\allowbreak{}ew}, rows \texttt{mt}, \texttt{als} \\
\ensuremath{\xi}\_\allowbreak{}inf(\ensuremath{\mu}\_\allowbreak{}b) 0.03030/\allowbreak{}0.03781/\allowbreak{}0.04413 & same : \texttt{xi\_\allowbreak{}inf\_\allowbreak{}mub} \\
\ensuremath{\alpha}\_\allowbreak{}s = 0.1171/\allowbreak{}0.1189 rows & same \\
MMR 2021 reproduction 0.05946 & \texttt{r3/\allowbreak{}inflation\_\allowbreak{}stability/\allowbreak{}run01/\allowbreak{}repro/\allowbreak{}repro\_\allowbreak{}step1.json}; \texttt{C155\_\allowbreak{}summary.json} : \texttt{reproduction} \\
all rows & \texttt{C155\_\allowbreak{}summary.json} : \texttt{xi\_\allowbreak{}inf\_\allowbreak{}rows} \\
\addlinespace
\multicolumn{2}{@{}>{\raggedright\arraybackslash}p{0.97\textwidth}@{}}{\textbf{\texttt{higgs/\allowbreak{}12\_\allowbreak{}two\_\allowbreak{}branch}} --- survival of both branches through the turn (D130 /\allowbreak{} C159) (Sec.~\ref{rep:h-survival})}\\*
P\ensuremath{_2} = P\ensuremath{_1}\ensuremath{^2} critical 172.376/\allowbreak{}172.759/\allowbreak{}173.124 (run02 base) & \texttt{r3/\allowbreak{}two\_\allowbreak{}branch/\allowbreak{}run01/\allowbreak{}C159\_\allowbreak{}summary.json} : \texttt{headline/\allowbreak{}crit\_\allowbreak{}as0p1180} \\
GH interference \ensuremath{\varepsilon} 0.9948--0.99997 & same : \texttt{headline/\allowbreak{}GH\_\allowbreak{}interference\_\allowbreak{}eps} \\
Fr\'echet bands [172.372, 172.498], [172.713, 172.954] & same : \texttt{crit} \\
static-spectrum check & \texttt{r3/\allowbreak{}two\_\allowbreak{}branch/\allowbreak{}run01/\allowbreak{}static\_\allowbreak{}check.json} : \texttt{runs} \\
\addlinespace
\multicolumn{2}{@{}>{\raggedright\arraybackslash}p{0.97\textwidth}@{}}{\textbf{\texttt{shape/\allowbreak{}01\_\allowbreak{}emptiness\_\allowbreak{}gate}} --- quanta budget, reheating coupling, latest preparation, photon-pair branch (Sec.~\ref{rep:gate})}\\*
\ensuremath{\Sigma}\textbar{}n\textbar{} \ensuremath{<} 4.762 (\ensuremath{\xi} = 0.2999), 5.047 (0.258), 4.898 (0.28), 6.80 (0); \ensuremath{\xi}(5) = 0.2649 (GATE-02..06) & \texttt{r3/\allowbreak{}shape\_\allowbreak{}hardening/\allowbreak{}run01/\allowbreak{}emptiness/\allowbreak{}C153\_\allowbreak{}emptiness.json} : \texttt{quanta\_\allowbreak{}budget/\allowbreak{}*/\allowbreak{}budget\_\allowbreak{}sum\_\allowbreak{}n}, \texttt{xi\_\allowbreak{}where\_\allowbreak{}budget\_\allowbreak{}equals/\allowbreak{}5} \\
AAA fermion budget 5.50 (GATE-07) & same : \texttt{quanta\_\allowbreak{}budget/\allowbreak{}*/\allowbreak{}fermion\_\allowbreak{}budget} \\
log\ensuremath{_1}\ensuremath{_0}R = 93.86 (free, j = 1); 94.04 (\ensuremath{\xi} = 0.30, j = 1); 96.48 (j = 85) (GATE-08..10) & \texttt{r5/\allowbreak{}C55\_\allowbreak{}window\_\allowbreak{}gate/\allowbreak{}C55\_\allowbreak{}window\_\allowbreak{}gate.json} : \texttt{slices/\allowbreak{}0/\allowbreak{}log10\_\allowbreak{}R}; \texttt{C153\_\allowbreak{}emptiness.json} : \texttt{reheating\_\allowbreak{}excess/\allowbreak{}xi0.30/\allowbreak{}log10R/\allowbreak{}\{1,85\}} \\
prefactor 10\textasciicircum{}97.79 (GATE-11) & \texttt{C55\_\allowbreak{}window\_\allowbreak{}gate.json} : \texttt{meta/\allowbreak{}log10\_\allowbreak{}prefactor} \\
T\_\allowbreak{}max = 7.264 /\allowbreak{} 7.177; 3-dof crossing 6.909 /\allowbreak{} 6.822 (GATE-12, 13) & \texttt{C153\_\allowbreak{}emptiness.json} : \texttt{preparation\_\allowbreak{}window/\allowbreak{}hLb0.0044/\allowbreak{}*}, \texttt{hLb0.0048/\allowbreak{}*} \\
side ratio 4.96 (\ensuremath{\xi} = 0.30), 1.74 (free), 3.56--5.43 over the \ensuremath{\xi} band (GATE-14, 16) & \texttt{fu/\allowbreak{}C80/\allowbreak{}C80\_\allowbreak{}results.json} : \texttt{runs/\allowbreak{}*} \\
\addlinespace
\multicolumn{2}{@{}>{\raggedright\arraybackslash}p{0.97\textwidth}@{}}{\textbf{\texttt{shape/\allowbreak{}02\_\allowbreak{}realtime\_\allowbreak{}theorem}} --- real-moment theorem, G(t\_\allowbreak{}x), closed-form floor and the conditional theorem t\_\allowbreak{}* (Sec.~\ref{rep:quantum-shape})}\\*
m\ensuremath{_1}\ensuremath{_1} = 10.554, m\ensuremath{_1}\ensuremath{_2} = 2.4685, m\ensuremath{_2}\ensuremath{_1} = 136.37, m\ensuremath{_2}\ensuremath{_2} = 31.99; K(19.3237) = 24.5630; K(18.0) = 20.3204 (FLOOR-01..03, 11) & \texttt{r3/\allowbreak{}shape\_\allowbreak{}hardening/\allowbreak{}run01/\allowbreak{}linear/\allowbreak{}C153\_\allowbreak{}linear.json} : \texttt{floor\_\allowbreak{}closed\_\allowbreak{}form/\allowbreak{}\{19.3237,18.0\}/\allowbreak{}*} \\
\ensuremath{\sigma}\_\allowbreak{}floor 0.10374 (H\_\allowbreak{}eff), 0.10403 (C60/\allowbreak{}C76), 0.1031 (earlier \ensuremath{\delta}); next order 0.09337 (FLOOR-04..08) & same : \texttt{floor\_\allowbreak{}conventions/\allowbreak{}*}, \texttt{next\_\allowbreak{}order\_\allowbreak{}floor/\allowbreak{}*} \\
G = 2.3952, 8.2632, 45.097, 313.79, 961.71; natural/\allowbreak{}floor 1.000485 (TX-01..06, FLOOR-10) & \texttt{r1/\allowbreak{}state\_\allowbreak{}selection/\allowbreak{}run03/\allowbreak{}C137c\_\allowbreak{}results.json} : \texttt{rows[k]/\allowbreak{}opt\_\allowbreak{}over\_\allowbreak{}floor}, \texttt{rows/\allowbreak{}0/\allowbreak{}adiabatic\_\allowbreak{}over\_\allowbreak{}floor} \\
G(\ensuremath{\infty}) = 2\ensuremath{\mathcal{G}} = 1348.55 (\ensuremath{\kappa} = 19.3237), 1058.34 (18.0) (TX-07, 08) & \texttt{C153\_\allowbreak{}linear.json} : \texttt{floor\_\allowbreak{}closed\_\allowbreak{}form/\allowbreak{}*/\allowbreak{}G\_\allowbreak{}inf\_\allowbreak{}closed\_\allowbreak{}2calG} \\
S\ensuremath{_2}(\ensuremath{\tau}) grid; t\_\allowbreak{}*(99 \%) = 0.23696 (MC 2.5--97.5 \%: 0.2245--0.2518); t\_\allowbreak{}*(90 \%) logistic 0.352; next order 0.26514 (TX-19..24) & \texttt{r3/\allowbreak{}shape\_\allowbreak{}hardening/\allowbreak{}run02/\allowbreak{}C153\_\allowbreak{}run02\_\allowbreak{}summary.json} : \texttt{lo\_\allowbreak{}t/\allowbreak{}t\_\allowbreak{}star\_\allowbreak{}linear\_\allowbreak{}interp/\allowbreak{}0.99/\allowbreak{}*}, \texttt{lo\_\allowbreak{}t/\allowbreak{}t\_\allowbreak{}star\_\allowbreak{}logistic/\allowbreak{}t\_\allowbreak{}star/\allowbreak{}0.9}, \texttt{nx\_\allowbreak{}t/\allowbreak{}...}; \texttt{r3/\allowbreak{}shape\_\allowbreak{}hardening/\allowbreak{}run01/\allowbreak{}C153\_\allowbreak{}nl\_\allowbreak{}summary.json} : \texttt{grid[]} \\
never-real two-branch survival 0.043 /\allowbreak{} 0.027; incoming single-branch 0.887 /\allowbreak{} 0.723 (TX-13..15) & \texttt{C153\_\allowbreak{}nl\_\allowbreak{}summary.json} : \texttt{never[*]} \\
adiabatic states: survival 0.992 /\allowbreak{} 0.609 /\allowbreak{} 0.129 (TX-18) & \texttt{C153\_\allowbreak{}nl\_\allowbreak{}summary.json} : \texttt{adiab[*]/\allowbreak{}future/\allowbreak{}p} \\
all of the above in one file & \texttt{r3/\allowbreak{}shape\_\allowbreak{}hardening/\allowbreak{}run01/\allowbreak{}C153\_\allowbreak{}summary.json} \\
\addlinespace
\multicolumn{2}{@{}>{\raggedright\arraybackslash}p{0.97\textwidth}@{}}{\textbf{\texttt{shape/\allowbreak{}03\_\allowbreak{}fossil}} --- the quantum fossil: floor, natural state, \ensuremath{\sigma}\_\allowbreak{}E, side difference and nonlinear ratio (Sec.~\ref{rep:quantum-shape})}\\*
floor, natural/\allowbreak{}floor, gain 216.46 (free) & \texttt{r5/\allowbreak{}C58\_\allowbreak{}turning\_\allowbreak{}point\_\allowbreak{}states/\allowbreak{}C58\_\allowbreak{}results.json} : \texttt{SM\_\allowbreak{}AAA/\allowbreak{}aspect/\allowbreak{}*} \\
\ensuremath{\sigma}\_\allowbreak{}E: 0.0589 (free), 0.1041 (\ensuremath{\xi} = 0.2999, C60/\allowbreak{}C76 convention), vs \ensuremath{\xi} (FOS-01, 02; FLOOR-06) & \texttt{fu/\allowbreak{}C82/\allowbreak{}C82\_\allowbreak{}results.json} : \texttt{sigma\_\allowbreak{}E\_\allowbreak{}leading/\allowbreak{}*} \\
side difference 18 \% /\allowbreak{} 10 \%; median coefficient 1.588 (FOS-04, 05) & same : \texttt{typical\_\allowbreak{}side\_\allowbreak{}difference\_\allowbreak{}plateau/\allowbreak{}*}, \texttt{median\_\allowbreak{}coeff} \\
plateau free \ensuremath{\sigma}\_\allowbreak{}E 0.0587 & \texttt{r5/\allowbreak{}C63\_\allowbreak{}rolling\_\allowbreak{}fossil/\allowbreak{}C63\_\allowbreak{}results.json} : \texttt{SM\_\allowbreak{}AAA/\allowbreak{}modes/\allowbreak{}aspect/\allowbreak{}natural\_\allowbreak{}rms} \\
K(19.32) = 24.5501; dln\ensuremath{\sigma}/\allowbreak{}dln\ensuremath{\kappa} = 1.366; LiteBIRD 3.1 \%, CMB-S4 1.6 \%; \ensuremath{\sigma} ratio 0.6306 at r = 1e-4 (FLOOR-01, FOS-14..16) & \texttt{r2/\allowbreak{}observables/\allowbreak{}run01/\allowbreak{}C143\_\allowbreak{}observables.json} : \texttt{F\_\allowbreak{}kappaE}, \texttt{dlnsigma\_\allowbreak{}dlnkappa}, \texttt{LiteBIRD\_\allowbreak{}*}, \texttt{CMBS4\_\allowbreak{}*}, \texttt{sigma\_\allowbreak{}vs\_\allowbreak{}r} \\
NL/\allowbreak{}linear aspect variance 1.043 /\allowbreak{} 1.045 (FOS-07) & \texttt{fu/\allowbreak{}C76/\allowbreak{}C76\_\allowbreak{}summary.json} : \texttt{sets/\allowbreak{}\{hi,ind\}/\allowbreak{}*} ; also \texttt{../\allowbreak{}04} C142 \texttt{regression/\allowbreak{}C76\_\allowbreak{}reference} \\
fossil distributions (free 2048, \ensuremath{\xi} = 0.30 2048) & \texttt{r5/\allowbreak{}nonlinear\_\allowbreak{}SM/\allowbreak{}C60\_\allowbreak{}summary.json}, \texttt{fu/\allowbreak{}C76/\allowbreak{}C76\_\allowbreak{}summary.json} + trajectory archives \\
\ensuremath{\phi}\_\allowbreak{}b, N\_\allowbreak{}post, hL\_\allowbreak{}b of the SM pivot & \texttt{r5/\allowbreak{}C62\_\allowbreak{}SM\_\allowbreak{}pivot/\allowbreak{}C62\_\allowbreak{}results.json} : \texttt{SM\_\allowbreak{}AAA/\allowbreak{}*} \\
\addlinespace
\multicolumn{2}{@{}>{\raggedright\arraybackslash}p{0.97\textwidth}@{}}{\textbf{\texttt{shape/\allowbreak{}04\_\allowbreak{}nonlinear\_\allowbreak{}survival}} --- survival thresholds, real states set in the past, plateau background, next order (Sec.~\ref{rep:quantum-shape})}\\*
survival vs s (NL-01) & \texttt{r2/\allowbreak{}shape\_\allowbreak{}nonlinear\_\allowbreak{}threshold/\allowbreak{}run01/\allowbreak{}C142\_\allowbreak{}summary.json} : \texttt{sets/\allowbreak{}amp\_\allowbreak{}s*/\allowbreak{}survival} \\
s\ensuremath{_5}\ensuremath{_0}, s\ensuremath{_9}\ensuremath{_0} (N = 256) & same : \texttt{amp\_\allowbreak{}threshold/\allowbreak{}*} \\
s\ensuremath{_9}\ensuremath{_9} = 2.756 (MC 95 \% 2.66--2.90), logistic 2.863 (boot 2.74--2.96); next order 3.160 /\allowbreak{} 3.294 (NL-06, 07, 09) & \texttt{r3/\allowbreak{}shape\_\allowbreak{}hardening/\allowbreak{}run02/\allowbreak{}C153\_\allowbreak{}run02\_\allowbreak{}summary.json} : \texttt{lo\_\allowbreak{}amp/\allowbreak{}s99\_\allowbreak{}loglinear/\allowbreak{}central}, \texttt{lo\_\allowbreak{}amp/\allowbreak{}s99\_\allowbreak{}logistic\_\allowbreak{}in\_\allowbreak{}ln\_\allowbreak{}s/\allowbreak{}s99}, \texttt{nx\_\allowbreak{}amp/\allowbreak{}*} \\
survival 0.9976 /\allowbreak{} 0.9834 /\allowbreak{} 0.9536 at s = 2.5 /\allowbreak{} 3.0 /\allowbreak{} 3.5 (N = 2048) (NL-03) & same : \texttt{lo\_\allowbreak{}amp/\allowbreak{}rows[*]/\allowbreak{}p} \\
past-set real states 0.623 /\allowbreak{} 0.184 /\allowbreak{} 0.133 (TX-27, 29) & \texttt{C142\_\allowbreak{}summary.json} : \texttt{sets/\allowbreak{}r1\_\allowbreak{}tx*/\allowbreak{}survival} \\
next order s\ensuremath{_9}\ensuremath{_0} = 4.19, s\ensuremath{_5}\ensuremath{_0} = 9.07 (NL-10) & \texttt{r3/\allowbreak{}shape\_\allowbreak{}hardening/\allowbreak{}run01/\allowbreak{}C153\_\allowbreak{}nl\_\allowbreak{}summary.json} : \texttt{next\_\allowbreak{}amp\_\allowbreak{}threshold} \\
plateau background: 0 /\allowbreak{} 2560 fate flips (NL-14) & \texttt{r3/\allowbreak{}shape\_\allowbreak{}tier2/\allowbreak{}run01/\allowbreak{}plateau/\allowbreak{}C160\_\allowbreak{}plateau\_\allowbreak{}summary.json} : \texttt{sets/\allowbreak{}*/\allowbreak{}n\_\allowbreak{}flips} \\
\addlinespace
\multicolumn{2}{@{}>{\raggedright\arraybackslash}p{0.97\textwidth}@{}}{\textbf{\texttt{shape/\allowbreak{}05\_\allowbreak{}moyal\_\allowbreak{}wdw}} --- SM Moyal correction, shear criterion with a quantum volume, B1 Wheeler--DeWitt test, B3 width lemma (Sec.~\ref{rep:quantum-shape})}\\*
Moyal \ensuremath{-}2.03e-4 /\allowbreak{} \ensuremath{-}1.08e-3 (\ensuremath{\xi} = 0.30); next order \ensuremath{-}1.79e-4 /\allowbreak{} \ensuremath{-}9.2e-4 (FOS-08..10) & \texttt{r3/\allowbreak{}shape\_\allowbreak{}tier2/\allowbreak{}run01/\allowbreak{}moyal/\allowbreak{}C160\_\allowbreak{}moyal.json} : \texttt{SM/\allowbreak{}xi030/\allowbreak{}rows/\allowbreak{}0/\allowbreak{}rel/\allowbreak{}a\_\allowbreak{}Heff/\allowbreak{}\{aspect,angle\}}, \texttt{SM/\allowbreak{}next\_\allowbreak{}order/\allowbreak{}...} \\
exact 1D quantum \ensuremath{-} TWA (FOS-13 analogue) & \texttt{.../\allowbreak{}moyal/\allowbreak{}C160\_\allowbreak{}moyal\_\allowbreak{}1d.json} : \texttt{checks[*]} \\
\ensuremath{\Delta}\ensuremath{\rho} = 7.8e-5 (NL-17) & \texttt{r3/\allowbreak{}shape\_\allowbreak{}tier2/\allowbreak{}run01/\allowbreak{}shear/\allowbreak{}C160\_\allowbreak{}shear\_\allowbreak{}wdw.json} : \texttt{physical}, \texttt{fits} \\
D\_\allowbreak{}re = 2.388e-6 at X = 2282.49; 3.1e-2 /\allowbreak{} 4.3e-4 /\allowbreak{} 1.5e-5 at X = 100 /\allowbreak{} 300 /\allowbreak{} 1000 (TS-01, 02) & \texttt{r3/\allowbreak{}shape\_\allowbreak{}tier2/\allowbreak{}run01/\allowbreak{}b1/\allowbreak{}b1\_\allowbreak{}cases.json} : \texttt{cases/\allowbreak{}A\_\allowbreak{}Ai\_\allowbreak{}real@<X>/\allowbreak{}rows/\allowbreak{}0/\allowbreak{}D\_\allowbreak{}re} \\
R15 counterexample D\_\allowbreak{}re 0.969 (\ensuremath{\varepsilon} = 1), 0.161 (\ensuremath{\varepsilon} = 0.3) (TS-03); (Ai+0.7Bi)\ensuremath{\otimes}\ensuremath{\chi}\ensuremath{_0} 5.0e-4 (TS-04) & \texttt{b1\_\allowbreak{}cases.json} /\allowbreak{} \texttt{.../\allowbreak{}b1/\allowbreak{}C160\_\allowbreak{}b1\_\allowbreak{}wdw\_\allowbreak{}zs0.json} : \texttt{runs[*]} \\
B3: Var H\_\allowbreak{}min = k/\allowbreak{}2m at width 1/\allowbreak{}(2m\textbar{}\ensuremath{\omega}\textbar{}) (TS-05) & stdout \texttt{logs/\allowbreak{}05\_\allowbreak{}moyal\_\allowbreak{}wdw/\allowbreak{}C141\_\allowbreak{}minvar\_\allowbreak{}lemma.out.txt} \\
all parts & \texttt{r3/\allowbreak{}shape\_\allowbreak{}tier2/\allowbreak{}run01/\allowbreak{}C160\_\allowbreak{}summary.json} \\
\addlinespace
\multicolumn{2}{@{}>{\raggedright\arraybackslash}p{0.97\textwidth}@{}}{\textbf{\texttt{shape/\allowbreak{}06\_\allowbreak{}primordial\_\allowbreak{}spectra}} --- n\_\allowbreak{}s and r of the plateau after the turn (Sec.~\ref{rep:spectra})}\\*
n\_\allowbreak{}s = 0.9653407, r = 0.0034624 (N\_\allowbreak{}* = 55.722) (PRIM-01, 02) & \texttt{r5/\allowbreak{}C64\_\allowbreak{}plateau\_\allowbreak{}ns\_\allowbreak{}r/\allowbreak{}C64\_\allowbreak{}results.json} : \texttt{rows/\allowbreak{}6/\allowbreak{}\{ns,r\}} \\
n\_\allowbreak{}s = 0.9617386, r = 0.0042022 (N\_\allowbreak{}* = 50.328) (PRIM-03, 04) & same : \texttt{rows/\allowbreak{}3/\allowbreak{}\{ns,r\}} \\
12 values of N\_\allowbreak{}* & same : \texttt{rows[*]} \\
initial-state tests & same : \texttt{initial\_\allowbreak{}state\_\allowbreak{}tests} \\
\addlinespace
\multicolumn{2}{@{}>{\raggedright\arraybackslash}p{0.97\textwidth}@{}}{\textbf{\texttt{cmb/\allowbreak{}01\_\allowbreak{}fisher\_\allowbreak{}kl}} --- torus T+E covariance, shape Fisher and exact KL (ideal full sky) (Sec.~\ref{rep:sky})}\\*
KL D\_\allowbreak{}TE, D\_\allowbreak{}T of the free fossil at 1.07 (end-to-end check) & \texttt{checks/\allowbreak{}kl\_\allowbreak{}quick.json} : \texttt{D\_\allowbreak{}TE\_\allowbreak{}fresh\_\allowbreak{}fresh}, \texttt{D\_\allowbreak{}T\_\allowbreak{}fresh\_\allowbreak{}fresh} \\
Fisher \ensuremath{\sigma}(\ensuremath{\beta}\_\allowbreak{}E) = 0.023 /\allowbreak{} 0.073 /\allowbreak{} 0.25 at L\ensuremath{_0} = 1.05 /\allowbreak{} 1.1 /\allowbreak{} 1.2 (CMB-23, earlier values) & \texttt{r5/\allowbreak{}obs\_\allowbreak{}fisher/\allowbreak{}C59\_\allowbreak{}shape\_\allowbreak{}fisher.json} (\texttt{--full}) \\
earlier window 1.093--1.133 for the typical fossil, circle bound 1.0931 (CMB-06, 22) & \texttt{fu/\allowbreak{}C75/\allowbreak{}C75\_\allowbreak{}summary.json} : \texttt{per\_\allowbreak{}sigma/\allowbreak{}*} \\
transfer-table metadata, low-k fix & \texttt{r5/\allowbreak{}obs\_\allowbreak{}fisher/\allowbreak{}C59\_\allowbreak{}transfer.json} \\
\addlinespace
\multicolumn{2}{@{}>{\raggedright\arraybackslash}p{0.97\textwidth}@{}}{\textbf{\texttt{cmb/\allowbreak{}02\_\allowbreak{}ensemble}} --- E-plane fossil ensemble, read edge and the P(read \textbar{} L\ensuremath{_0}) window (C145/\allowbreak{}D117) (Sec.~\ref{rep:sky})}\\*
P\_\allowbreak{}joint peak 0.3440 at 1.075 (CMB-02) & \texttt{ai/\allowbreak{}cmb\_\allowbreak{}ensemble/\allowbreak{}run01/\allowbreak{}results/\allowbreak{}S3\_\allowbreak{}postcheck.json} : \texttt{variants/\allowbreak{}TE80\_\allowbreak{}grid/\allowbreak{}modes/\allowbreak{}mid/\allowbreak{}\{Pj\_\allowbreak{}max,L0\_\allowbreak{}Pj\_\allowbreak{}max\}} \\
P\_\allowbreak{}joint \ensuremath{\ge} 0.1 for [1.0175, 1.1825]; \ensuremath{\ge} 0.3 /\allowbreak{} 0.2 /\allowbreak{} 0.05 windows (CMB-03, 04) & same : \texttt{.../\allowbreak{}win\_\allowbreak{}joint/\allowbreak{}\{0.1,0.3,0.2,0.05\}} \\
integrated measure 0.04223 (CMB-05); representative window [1.0931, 1.1394] (CMB-06) & same : \texttt{.../\allowbreak{}mean\_\allowbreak{}span}, \texttt{.../\allowbreak{}rep\_\allowbreak{}window} \\
P\_\allowbreak{}cond \ensuremath{\ge} 0.5 up to 1.090 (CMB-07) & same : \texttt{.../\allowbreak{}win\_\allowbreak{}cond/\allowbreak{}0.5} \\
read edge u\_\allowbreak{}max 1.021--1.039 (CMB-01) & same : \texttt{variants/\allowbreak{}TE80\_\allowbreak{}grid/\allowbreak{}edges} \\
noise 0.313 /\allowbreak{} 0.0384; T only 0.286 /\allowbreak{} 0.0352 (CMB-10, 11) & same : \texttt{variants/\allowbreak{}\{noise\_\allowbreak{}TE80\_\allowbreak{}grid,T80\_\allowbreak{}grid\}/\allowbreak{}modes/\allowbreak{}mid/\allowbreak{}*} \\
next order 0.377 at 1.065; [1.015, 1.1675]; 0.0421; rep window [1.082, 1.128] (CMB-14..17) & \texttt{r3/\allowbreak{}shape\_\allowbreak{}hardening/\allowbreak{}run01/\allowbreak{}ensemble/\allowbreak{}C153\_\allowbreak{}ensemble.json} : \texttt{results/\allowbreak{}0.0940/\allowbreak{}variants/\allowbreak{}TE80\_\allowbreak{}grid/\allowbreak{}mid/\allowbreak{}*} \\
55 fossils (columns: xi = sin 3\ensuremath{\theta}, s = L\ensuremath{_0}/\allowbreak{}d\_\allowbreak{}LSS, u) & \texttt{r3/\allowbreak{}cleanup\_\allowbreak{}t12/\allowbreak{}run01/\allowbreak{}fossil\_\allowbreak{}table\_\allowbreak{}v2.csv} \\
near-window packet 1.7 \%, \ensuremath{\alpha} = 17.6 for P = 0.5 (CMB-12, 13) & \texttt{ai/\allowbreak{}cmb\_\allowbreak{}ensemble/\allowbreak{}run01/\allowbreak{}priors/\allowbreak{}priors\_\allowbreak{}summary.json}, \texttt{P4\_\allowbreak{}size\_\allowbreak{}prior\_\allowbreak{}ens.json} \\
KL table, profiles (TE30/\allowbreak{}40/\allowbreak{}80, T80, fixed/\allowbreak{}profiled), Brent, convergence & \texttt{ai/\allowbreak{}cmb\_\allowbreak{}ensemble/\allowbreak{}run01/\allowbreak{}results/\allowbreak{}\{kl\_\allowbreak{}table,S2\_\allowbreak{}profiles,S2\_\allowbreak{}brent,S4\_\allowbreak{}convergence,S4\_\allowbreak{}noise,S3\_\allowbreak{}analysis\}.json} (archived; \texttt{--full} recomputes) \\
\addlinespace
\multicolumn{2}{@{}>{\raggedright\arraybackslash}p{0.97\textwidth}@{}}{\textbf{\texttt{cmb/\allowbreak{}03\_\allowbreak{}joint\_\allowbreak{}state}} --- joint inflaton--shape state and the size of the torus (C147/\allowbreak{}D119) (Sec.~\ref{rep:sky})}\\*
\ensuremath{\sigma}\_\allowbreak{}\ensuremath{\phi} = 0.00973692 (JOINT-01) & \texttt{ai/\allowbreak{}joint\_\allowbreak{}state/\allowbreak{}run01/\allowbreak{}joint\_\allowbreak{}state\_\allowbreak{}results.json} : \texttt{T3/\allowbreak{}B3/\allowbreak{}ref/\allowbreak{}sigma\_\allowbreak{}B3} \\
\ensuremath{\sigma}\_\allowbreak{}N \ensuremath{\ge} 0.935075 (ref) /\allowbreak{} 0.938387 (window) (JOINT-02) & same : \texttt{T3/\allowbreak{}response/\allowbreak{}\{ref,window\}/\allowbreak{}floor\_\allowbreak{}sigma\_\allowbreak{}N} \\
\ensuremath{\Delta}\ensuremath{\phi}\_\allowbreak{}b = 6.1825e-4; N window [68.480, 68.516], width 0.0359 (JOINT-04, 05) & same : \texttt{T4/\allowbreak{}typical\_\allowbreak{}fossil\_\allowbreak{}window/\allowbreak{}dphi\_\allowbreak{}window}, \texttt{T2/\allowbreak{}N\_\allowbreak{}window\_\allowbreak{}typical}, \texttt{T2/\allowbreak{}window\_\allowbreak{}width\_\allowbreak{}efold} \\
flat-d\ensuremath{\phi} readability 1.389e-3 (N\_\allowbreak{}max = 100) \ensuremath{\to} 5.35e-5 (1e6) (JOINT-06) & same : \texttt{T4/\allowbreak{}shape\_\allowbreak{}marginalised/\allowbreak{}A\_\allowbreak{}fixed\_\allowbreak{}typical/\allowbreak{}flat\_\allowbreak{}dphi\_\allowbreak{}Nmax=*} \\
S\_\allowbreak{}E = 11964 (JOINT-09); \ensuremath{\phi}\_\allowbreak{}b(ref) = 5.57936 (BKG-03) & same : \texttt{T2/\allowbreak{}S\_\allowbreak{}E\_\allowbreak{}ref\_\allowbreak{}closed}, \texttt{T2/\allowbreak{}phib\_\allowbreak{}ref(N=68.25)} \\
interface to C145 & \texttt{ai/\allowbreak{}joint\_\allowbreak{}state/\allowbreak{}run01/\allowbreak{}C145\_\allowbreak{}interface.json} \\
symbolic slope identities & stdout of \texttt{b3\_\allowbreak{}slope\_\allowbreak{}symbolic.py} \\
\addlinespace
\multicolumn{2}{@{}>{\raggedright\arraybackslash}p{0.97\textwidth}@{}}{\textbf{\texttt{validation}} --- independent reimplementations (blind keystone and V1--V5, semi-blind T19, compute-first R14 /\allowbreak{} R15) and their comparison with the main calculation; archived files relative to \texttt{validation/} (Sec.~\ref{rep:v-blind})}\\*
23 compared numbers (keystone 10, V2 2, V1 3, V3 2, V4 4, V5 2), all within the pass criteria & \texttt{output/\allowbreak{}validation/\allowbreak{}comparison.json} : \texttt{items}; \texttt{COMPARISON.md} (\texttt{compare\_\allowbreak{}blind.py}) \\
keystone (Table~\ref{tab:v19-keystone}) and V2 cross terms & \texttt{keystone/\allowbreak{}results/\allowbreak{}BLIND\_\allowbreak{}RESULTS.json}, \texttt{keystone/\allowbreak{}results\_\allowbreak{}v2/\allowbreak{}BLIND\_\allowbreak{}RESULTS\_\allowbreak{}V2.json} \\
V1 SM mapping; V3 claims; V4 shape; V5 universality and inflation & \texttt{sm\_\allowbreak{}mapping/}, \texttt{claims/}, \texttt{shape/}, \texttt{universality\_\allowbreak{}inflation/} (each \texttt{SPEC.md}, \texttt{code/}, \texttt{results/}) \\
T19 (semi-blind); R14, R15 (compute first) & \texttt{T19\_\allowbreak{}semiblind/}, \texttt{early\_\allowbreak{}rounds/\allowbreak{}R14/}, \texttt{early\_\allowbreak{}rounds/\allowbreak{}R15/} \\
intermediate steps of report Secs.~1--5 and 13--17 & \texttt{report\_\allowbreak{}checks/\allowbreak{}rep-A\_\allowbreak{}checks.py}, \texttt{rep-A\_\allowbreak{}checks\_\allowbreak{}part2.py} (rerun by \texttt{run.py} against \texttt{expected/\allowbreak{}*.out.txt}) \\
\addlinespace
\end{longtable}}


\subsection{What the package does not contain}\label{sec:v22-limits}

Large intermediate files are not shipped but are regenerated by the full runs: the covariance caches of the microwave-background chains (1.3~GB, 0.7~GB, 3.3~GB and 15.5~GB) and the Higgs node tables and gauge profiles (40--170~MB per run).
The transfer tables for $\ell\le40$ (34~MB) and $\ell\le80$ (85~MB) and the archives of the nonlinear trajectories (69~MB in six files) are shipped; the table for $\ell\le80$, needed only for \texttt{cmb/02 --full}, was computed with CAMB~2.0.4 by a script of the separate line of calculation of Sec.~\ref{rep:v-conditional}; the script and its settings are shipped with the table in \texttt{cmb/\allowbreak 02\_ensemble/\allowbreak data/}. %
The blind reimplementations of Sec.~\ref{rep:v-blind} are described there and archived in the package's validation folder, \texttt{validation/}, whose script compares $23$ of their numbers with the main calculation; all agree within the pass criteria of the specifications. The blind recalculation of the abelian model (Sec.~\ref{sec:v11-blind-abelian}) is archived from version~1.2.0 in \texttt{validation/\allowbreak abelian\_blind/}, with its comparison in \texttt{post\_comparison/\allowbreak COMPARISON.md}. %
The original runs of the first two stages of the Higgs calculation are not repeated; where a later step depends on them, their archived outputs are staged as inputs.
The $66$ node jobs that split the cross-term shift by groups of fields (Sec.~\ref{sec:rb-xfinal}) are archived and not rerun; \texttt{higgs/\allowbreak 04\_cross\_terms/\allowbreak cross\_decomp.py} recomputes the split from them in both modes. %
Each chain README has a section ``Not packaged / partial coverage'' that lists the remaining gaps, for example the original scripts of the run02-based extras of the real-time step, whose outputs are archived.
Some archived JSON and Markdown files contain absolute paths of the research machine in metadata fields; they are kept unchanged so that their SHA-256 matches the research archive, and the code never uses these fields.

\subsection{The static-box and abelian calculations}\label{sec:v11-repro-supplement}

The folder \texttt{prl\_supplement/} holds the Supplemental Material of the Letter, unchanged: the scripts and archived outputs of the static box (Sec.~\ref{sec:v11-static-box}; folders \texttt{c170/} and \texttt{results/C170/}) and of the abelian Higgs model (Sec.~\ref{sec:v11-abelian}; \texttt{c171/} and \texttt{results/C171/}), the script of Fig.~\ref{fig:v11-prl-fig1} (\texttt{figures/}), the generator of \texttt{numbers\_c170.tex} (\texttt{numbers/}), its own \texttt{README.txt} and \texttt{SHA256SUMS.txt}, and a note on its use inside the package (\texttt{PACKAGE\_NOTE.md}).
The point-by-point records of the two calculations are bundled in two JSON-Lines files, which \texttt{unbundle\_records.py} restores and checks against their SHA-256.
The static-box scripts import the production modules of the Higgs chain (\texttt{higgs/\allowbreak zmlib}, \texttt{higgs/\allowbreak 04\_cross\_terms}, \texttt{dnec\_common}), which are unchanged since version~1.0.0; \texttt{c170/\allowbreak setup\_from\_release.py} copies them from the folder named by the environment variable \texttt{DNEC\_RELEASE}, inside the package its root, and verifies their SHA-256.
The abelian solver needs only NumPy and SciPy.

The README gives three routes: (A) the analysis from the archived point records, in minutes; (B) the full recomputation, several CPU hours for the static box and about $7$~min for the abelian model; (C) the figure and the macros. %
Route~A with \texttt{DNEC\_RELEASE} set to version~1.1.0 of the package verifies the SHA-256 of all copied modules, reproduces the archived widths of the final model exactly (and to $7\times10^{-12}$ from generic starting points), passes all 30 checks of \texttt{results/\allowbreak C170/\allowbreak checks.json}, and reproduces the table of widths and the thresholds of the abelian model. %
Table~\ref{tab:v11-supplement-map} gives the output that contains each number of Secs.~\ref{sec:v11-static-box} and~\ref{sec:v11-abelian}; the README of the folder lists the same correspondence for the numbers of the Letter.
The numbers of Secs.~\ref{sec:v11-mechanism-context}, \ref{sec:v11-static-box}--\ref{sec:v11-static-torus} and~\ref{sec:v11-abelian} that are computed from these outputs, among them the WKB exponents and Fubini actions at the thresholds, the closed-form deviations, the local exponent of the widths, the $\kappa$ ratios, the box ratio, the one-dimensional WKB decomposition of the slope and the abelian Born--Oppenheimer bracket, are produced by \texttt{figures/\allowbreak derive\_rep\_v11\_static\_box.py}, \texttt{figures/\allowbreak derive\_rep\_v11\_wkb\_split.py} and \texttt{figures/\allowbreak tab\_rep\_v11\_abelian.py} of the package, and the figures of these sections other than Fig.~\ref{fig:v11-prl-fig1} by the scripts \texttt{figures/\allowbreak fig\_rep\_v11\_*.py}; the inputs of these scripts are copies of the archived outputs in \texttt{figures/\allowbreak data/\allowbreak higgs/\allowbreak static\_box\_c170/} and \texttt{figures/\allowbreak data/\allowbreak higgs/\allowbreak abelian\_c171/}.

\begin{table}[ht]
\centering
\small
\caption{Numbers of the static box and of the abelian model and the archived outputs that contain them (paths under \texttt{prl\_supplement/\allowbreak results/} unless they start with \texttt{higgs/} or \texttt{figures/}; ``\#'' points to a key).}
\label{tab:v11-supplement-map}
\begin{tabular}{@{}p{0.43\textwidth}p{0.51\textwidth}@{}}
\toprule
Quantity & Output \\
\midrule
Width law in the fit window: $a$, $b$, residual, effective $\sigma$ & \texttt{C170/\allowbreak S1/\allowbreak C170\_S1\_summary.json\#boxes.A.fits} \\
Level below the barrier top, WKB exponents in the window & \texttt{C170/\allowbreak S1/\allowbreak C170\_S1\_summary.json\#boxes.A.table} \\
Zero-point slope and its reduction by $\Delta V$ & \texttt{C170/\allowbreak S2/\allowbreak C170\_S2\_curves.json\#slopes} \\
$\lamc(N)$ of boxes A--E, closed-form accuracy, $\lamc(300)$ across boxes & \texttt{C170/\allowbreak S1/\allowbreak C170\_S1\_lambda\_c\_N.csv} \\
Half-survival durations $N^*$ of the Standard Model vacuum & \texttt{C170/\allowbreak S1/\allowbreak C170\_S1\_lambdaSM\_Nstar.csv} \\
$N_{\rm eff}$; range of the narrow resonance & \texttt{C170/\allowbreak S1/\allowbreak C170\_S1\_summary.json\#boxes.A} \\
Regression; static minus hazard threshold & \texttt{C170/\allowbreak C170\_results.json\#S0} \\
Model error of $\lamc(300)$ & \texttt{higgs/\allowbreak 09\_claims/\allowbreak expected/\allowbreak w/\allowbreak claims\_run05/\allowbreak run01/\allowbreak table1\_rebuild.json\#composite/\allowbreak minus\_no\_vi/\allowbreak 0.5} (shallow side) and \texttt{plus\_no\_vi/\allowbreak 0.5} of the package, times $0.00659\GeV^{-1}$ \\
Cross-term shift of $\lamc(300)$ & \texttt{C170/\allowbreak C170\_results.json\#\allowbreak S2.\allowbreak kappa1\_vs\_\allowbreak final\_model} \\
$\kappa$ scan: thresholds, $\kappa_{\rm thr}$, $\kappa_{\rm min}$, potentials & \texttt{C170/\allowbreak S2/\allowbreak C170\_S2\_lambda\_c\_kappa.csv}, \texttt{C170\_results.json\#S2} \\
Coupling scans & \texttt{C170/\allowbreak S3/\allowbreak C170\_S3\_coupling\_scan.csv} \\
Abelian model: widths, local exponent, method cross-check & \texttt{C171/\allowbreak C171\_gamma\_table.csv}, \texttt{C171/\allowbreak C171\_cap\_crosscheck.csv} \\
Abelian model: thresholds, Born--Oppenheimer widths & \texttt{C171/\allowbreak C171\_thresholds.csv}, \texttt{C171/\allowbreak C171\_results.json} \\
Abelian model: holonomy tails & \texttt{C171/\allowbreak C171\_holonomy\_tails.csv} \\
Derived numbers of Secs.~\ref{sec:v11-mechanism-context}, \ref{sec:v11-static-box}--\ref{sec:v11-static-torus}: WKB exponents, Fubini actions, Planck operator, closed form, deep-well exponent, $\kappa$ ratios, box ratio & \texttt{figures/\allowbreak out/\allowbreak derived\_rep\_v11\_static\_box.json} (\texttt{figures/\allowbreak derive\_rep\_v11\_static\_box.py}) \\
WKB decomposition of the width slope; local exponent of the widths of box~A & \texttt{figures/\allowbreak out/\allowbreak derived\_rep\_v11\_wkb\_split.json} (\texttt{figures/\allowbreak derive\_rep\_v11\_wkb\_split.py}) \\
Derived numbers of Sec.~\ref{sec:v11-abelian}: barrier top, energy coefficient, Born--Oppenheimer bracket, Kemble widths, barrier at the Standard Model slope & \texttt{figures/\allowbreak tab\_rep\_v11\_abelian.py} \\
Figures~\ref{fig:v11-static}, \ref{fig:v11-kappa}, \ref{fig:v11-ab-width} and~\ref{fig:v11-ab-threshold} & \texttt{figures/\allowbreak fig\_rep\_v11\_*.py}, with inputs in \texttt{figures/\allowbreak data/\allowbreak higgs/\allowbreak static\_box\_c170/} and \texttt{abelian\_c171/} \\
\bottomrule
\end{tabular}
\end{table}
